% !TeX program = pdflatex
\pdfmapfile{+t1-times.map}
\pdfmapfile{+t1-formata.map}
\pdfmapfile{+t1-giovannistd.map}
\documentclass{ieeeaccess}
\usepackage{cite}
\usepackage{amsmath,amssymb,amsfonts}
\usepackage{graphicx}
\graphicspath{{../pic/}}
\usepackage{textcomp}
\usepackage{bm}
\usepackage{array}
\usepackage{booktabs}
\usepackage{placeins}
\usepackage{balance}
\setcounter{topnumber}{3}
\setcounter{dbltopnumber}{3}
\setcounter{totalnumber}{5}
\renewcommand{\topfraction}{0.95}
\renewcommand{\dbltopfraction}{0.95}
\renewcommand{\textfraction}{0.05}
\renewcommand{\floatpagefraction}{0.9}
\renewcommand{\dblfloatpagefraction}{0.9}
\usepackage{float}
\usepackage{algorithm}
\usepackage{algpseudocode}

\newcommand{\citep}[1]{\cite{#1}}

\def\BibTeX{{\rm B\kern-.05em{\sc i\kern-.025em b}\kern-.08em
    T\kern-.1667em\lower.7ex\hbox{E}\kern-.125emX}}

\begin{document}
\history{Date of publication xxxx 00, 0000, date of current version xxxx 00, 0000.}
\doi{10.1109/ACCESS.2026.DOI}

\title{Slope-Aware LPV-MPC Path Tracking of a Dual-Steering-Wheel AGV with Time-Lag-Enhanced ModernTCN Perception}
\author{\uppercase{Author Name}\authorrefmark{1}}
\address[1]{Affiliation placeholder.}
\tfootnote{This work was supported in part by xxxx under Grant xxxx.}

\markboth
{Author \headeretal: Slope-Aware LPV-MPC Path Tracking of a Dual-Steering-Wheel AGV}
{Author \headeretal: Slope-Aware LPV-MPC Path Tracking of a Dual-Steering-Wheel AGV}

\corresp{Corresponding author: Author Name (e-mail: email@example.com).}

\begin{abstract}

This paper proposes a slope-aware path-tracking method based on temporal-perception-enhanced linear parameter-varying model predictive control (LPV-MPC) to address the degradation of control accuracy and stability caused by omitting or inaccurately estimating the longitudinal road-grade angle in conventional path-tracking approaches. Firstly, to reduce reliance on potentially biased IMU-derived grade measurements, a lightweight Modern Temporal Convolutional Network (ModernTCN) is introduced to estimate the grade angle from historical onboard responses. The multi-scale time-lag offsets and the historical-window length are configured according to the vehicle response characteristics, yielding better adaptation to the operating conditions considered here than the original ModernTCN. Secondly, an uncertainty-adaptive ModernTCN--IMU interface treats the ModernTCN estimate as the primary source and uses the quality-qualified IMU estimate only as a bounded correction, thereby improving the reliability of the fused grade estimate under varying operating conditions. Furthermore, the conditioned fused-grade signal is synchronously incorporated into the LPV prediction model, measured-disturbance channel, grade-feedforward term, scheduled weights, and constraints. Compared with conventional LPV-MPC schemes that disregard grade information, this integration substantially enhances control precision and adaptation capability in graded-road scenarios. Experimental results demonstrate that the proposed method yields consistent closed-loop improvements in path-tracking accuracy and control-input regularity under longitudinal grade variations while retaining a lightweight computational structure.
\end{abstract}

\begin{keywords}

Dual-steering-wheel AGV, LPV-MPC, ModernTCN, path tracking, road-grade estimation

\end{keywords}

\titlepgskip=-21pt
\maketitle

\section{Introduction}

\PARstart{A}{utomated} guided vehicles (AGVs) are widely used in smart manufacturing, automated warehousing, and industrial logistics, where they undertake material handling and flexible transport tasks \citep{Zhang2023AGVSurvey,Oyekanlu2020AGVReview,Bhargava2024AGVReview}. Reliable AGV operation depends on several key technologies, including perception, localization, path planning, scheduling, and motion control. Among them, path tracking is the final execution layer that determines whether the planned route can be followed accurately, smoothly, and safely under actuator constraints and varying operating conditions. Therefore, robust path tracking control is a fundamental requirement for industrial AGV deployment \citep{DeRyck2020AGVControlReview}. The present study focuses on path tracking of a diagonal dual-steering-wheel (dual-drive-and-steer) AGV operating along industrial routes with longitudinal grade changes.

Linear parameter-varying model predictive control (LPV-MPC) has been widely adopted for vehicle and AGV path tracking because it handles multivariable dynamics, reference preview, and actuator constraints within a unified optimization framework \citep{Manikandan2022CurveAwareMPCAGV,Alcala2020AutonomousRacingLPVMPC,Kedir2024PSOLPVMPC,Domina2022LTVMPC}. In the considered dual-steering-wheel AGV, the vehicle speed, yaw rate, and longitudinal road-grade angle form the scheduling vector. The same grade angle is also used in the measured-disturbance channel, grade-feedforward term, scheduled weights, and input bounds. Uncompensated grade variation otherwise acts as a persistent disturbance that has motivated dedicated slope-aware assistance designs \citep{Luo2015DownhillAssist}. Omitting or inaccurately estimating this angle may therefore cause the LPV prediction model and optimization problem to be scheduled at an operating condition inconsistent with the actual plant, resulting in prediction mismatch and degraded closed-loop tracking and control performance \citep{Sato2013InexactScheduling,Pan2022InexactSchedulingLPV,Metzler2021PredictionFidelity,Lin2019AdaptiveMPCPathTracking}. A straightforward approach is to estimate the grade angle using an inertial measurement unit (IMU). Yet, IMU-derived grade estimates can be biased by transient acceleration, pitch dynamics, low-cost sensor noise, and road-induced vibrations \citep{Jauch2018IMUGradeOrientation,Liu2023RoadSlopeAAIMM,Gao2024RoadSlopeHDV}. While model-based observers, such as recursive least-squares and Kalman-filter variants, can infer road grade from vehicle dynamics, wheel-speed, and driving-force data \citep{Jiang2019AdaptiveSlopeMass,Guo2022SageHusaAKF}, their accuracy depends strongly on nominal vehicle parameters, load variations, dynamic maneuvers, and the validity of wheel-ground contact assumptions \citep{Gao2024RoadSlopeHDV}. These limitations motivate a data-driven temporal perception strategy. Unlike rigid physics-based observers, learning-based models can extract grade-related patterns from historical proprioceptive sequences. The coupled evolution of velocity attenuation, acceleration variation, yaw response, and drive current may therefore provide complementary information for LPV-MPC scheduling \citep{Liang2022LSTMRoadID}.

Within the realm of learning-based sequence modeling, one-dimensional convolutional neural networks (1D-CNNs) have been extensively adopted for engineering time-series analysis because of their efficiency in extracting local temporal patterns from sensor data \citep{Bueno2023Inception1DCNN}. However, when non-causal convolutional structures or offline-centered windows are directly used for online prediction, the feature representation at the current time step may be influenced by future samples, violating the causality requirement of real-time control systems. Temporal convolutional networks (TCNs) address this issue by enforcing causal convolutions, ensuring that each prediction depends only on historical and current inputs. By integrating causal structures with dilated convolutions and residual connections, TCNs effectively enlarge the receptive field and have been successfully applied to industrial sequence modeling tasks such as machinery remaining-useful-life prediction \citep{Bai2018TCN,Cao2021TCNRSA}.

Nevertheless, standard TCNs are not fully suited to industrial multivariate AGV perception. Their causal dilated convolutions enlarge the receptive field through fixed sampling intervals, but a mismatched dilation pattern may weaken short grade-sensitive transients \citep{Cao2021TCNRSA}. Meanwhile, AGV grade perception depends on coupled variables including velocity, acceleration, yaw response, and drive current, making cross-variable interaction important for reliable estimation \citep{Jin2023MTGODE}. ModernTCN revisits pure convolutional time-series modeling with large-kernel depthwise convolution for temporal dependency modeling and pointwise channel mixing for cross-variable interaction \citep{Luo2024ModernTCN}. These principles provide a suitable basis for the history-only lightweight estimator developed in this work.

However, constructing a lightweight history-based estimator from the temporal modeling principles of ModernTCN does not by itself resolve slope-aware path tracking. Its temporal input must preserve rapid grade-sensitive vehicle responses while providing sufficient historical context to distinguish persistent road inclination from transient acceleration, vibration, and local surface irregularities. The lag interval and historical window length therefore become distinct but coupled design choices. The problem is further complicated by the condition-dependent reliability of the learning-based and inertial estimates, as well as by the need to translate the estimated grade into control-relevant scheduling information for LPV-MPC. Learning-enhanced MPC for mobile platforms typically embeds learned components in the prediction model or controller \citep{Zhang2022LBMPCSurvey,Li2024PINNAGV}; here, learning stays on the perception side. These unresolved perception-to-control issues define the scope of the contributions presented below.

The methodological contributions are threefold:
\begin{enumerate}
\item Compared with direct IMU-based grade estimation
\citep{Jauch2018IMUGradeOrientation,Li2022QuaternionUKFGrade,
Kim2023InertialRoadPitch}, we adapt the established temporal-modeling
principles of ModernTCN \citep{Luo2024ModernTCN} to infer road grade
from history-only AGV responses without a direct grade input. To
improve grade-estimation accuracy, the task-specific estimator combines
a dynamics-informed 22-dimensional proprioceptive feature set with
multi-scale temporal differences tailored to AGV dynamics and
grade-signal characteristics. This representation retains
operating-condition context while capturing short grade-sensitive
responses, supporting more accurate grade estimation and improved
closed-loop control-input regularity.
\item Compared with using our adapted ModernTCN estimator alone, whose
temporal-modeling backbone is simplified from the established ModernTCN
architecture \citep{Luo2024ModernTCN}, we develop uncertainty-adaptive fusion that
uses a qualified causal RKF-IMU observer \citep{Chen2024RobustAttitude} only as a calibrated,
innovation-conditioned, bounded correction with exact fallback,
yielding lower mean scheduled-grade error and supported improvements
in heading tracking and control-input regularity.

\item Compared with LPV-MPC path tracking that omits grade scheduling
\citep{Alcala2020AutonomousRacingLPVMPC,
Tian2022GainScheduledLPVMPC,Cheng2021VelocityVaryingMPC}, we use one
conditioned fused-grade signal to synchronously update the prediction
model, measured-disturbance channel, grade feedforward, scheduled
weights, and constraints, enabling consistent grade-aware adaptation
while retaining model-based interpretability and constraint handling.
\end{enumerate}

The remainder of this paper is organized as follows.
Section~\ref{sec:agv_slope_estimation} establishes the model of the diagonal
dual-steering-wheel AGV and formulates the slope-aware path-tracking and
history-based grade-estimation problem.
Section~\ref{sec:proposed_framework} develops the proposed framework,
comprising the lightweight ModernTCN grade estimator, the dynamics-informed
feature and multi-scale delta-lag representation, the uncertainty-adaptive
fusion interface between ModernTCN and the qualified observer, and the fused-grade-scheduled LPV-MPC.
Section~\ref{sec:experiment_design} describes the simulation platform, the
frozen experimental configuration, and the evaluation protocol.
Section~\ref{sec:results_discussion} reports and discusses the offline and
closed-loop results. Finally, Section~\ref{sec:conclusion} concludes the
paper.

\section{AGV Modeling and History-Based Grade Estimation}

\label{sec:agv_slope_estimation}


\subsection{AGV Model}

\begin{figure}[!tb]
\centering
\includegraphics[width=\columnwidth]{fig02_agv_model.pdf}
\caption{Configuration of the diagonal dual-steering-wheel AGV. The LF and RR wheels are active drive--steer wheels, while the RF and LR wheels are passive support wheels. The principal motion variables and road-grade angle are indicated.}
\label{fig:agv_configuration}
\end{figure}


The considered diagonal dual-steering-wheel AGV has active drive--steer wheels at the left-front (LF) and right-rear (RR) positions, while the right-front (RF) and left-rear (LR) wheels provide passive vertical support. The normal loads of the passive wheels are retained in the load distribution and rolling resistance, whereas their lateral-force and yaw-moment contributions are neglected. The vehicle configuration and the principal motion variables are illustrated in Fig.~\ref{fig:agv_configuration}. Following the geometry-first modeling sequence used for multi-steering AGVs \citep{Li2025FourSteeringAGV,Wan2025AGVLSTMMPC}, the steering geometry is established before the actuator, load, tire-force, and body-dynamics models.

The vehicle motion-state and control vectors are
\begin{equation}
\begin{aligned}
\bm{x}&=
\begin{bmatrix}
X & Y & \psi & v & \omega & \delta_{lf} & \delta_{rr} & \beta
\end{bmatrix}^{T},\\
\bm{u}&=
\begin{bmatrix}
F_{cmd} & \omega_{cmd}
\end{bmatrix}^{T},
\end{aligned}
\label{eq:agv_state_control}
\end{equation}
where $(X,Y)$ and $\psi$ denote the position and heading in the two-dimensional road-aligned path coordinates, $v$ is the longitudinal speed, $\omega$ is the yaw rate, $\delta_{lf}$ and $\delta_{rr}$ are the active-wheel steering angles, and $\beta$ is the vehicle sideslip angle. The longitudinal road-grade angle $\theta$ is attached to the corresponding path station and enters the plant as an exogenous variable, with $\theta>0$ denoting an uphill grade along the positive direction of travel.

\subsubsection{Kinematic Model}

This study considers planar path tracking under longitudinal road grades. The reference route is represented in a two-dimensional road-aligned coordinate chart, while the grade angle is treated as a path-dependent exogenous variable. Vertical displacement, body-pitch dynamics, road bank, and cross-slope effects are outside the present scope. Consequently, the grade angle does not directly enter the planar kinematic mapping; instead, it affects the vehicle dynamics through grade resistance, grade-dependent normal loads, and tire-force limits. The planar vehicle kinematics are therefore
\begin{equation}
\dot X=v\cos(\psi+\beta),\qquad
\dot Y=v\sin(\psi+\beta),\qquad
\dot\psi=\omega.
\label{eq:agv_kinematics}
\end{equation}
Here, the longitudinal grade is assumed to be locally aligned with the path tangent.

For curve motion, the instantaneous center of rotation (ICR) is assumed to lie on the body $y$-axis. For the smoothed speed--yaw-rate pair $(v_{icr},\omega_{icr})$, its signed radius is $R_{icr}=v_{icr}\operatorname{sgn}(\omega_{icr})/\max(|\omega_{icr}|,\varepsilon)$, where $\varepsilon=10^{-6}$. With the center of mass midway between the axles, the active-wheel steering targets are
\begin{equation}
\begin{aligned}
\delta_{lf}^{tgt}&=\arctan\!\left(\frac{L_w/2}{R_{icr}-W/2}\right),\\
\delta_{rr}^{tgt}&=-\arctan\!\left(\frac{L_w/2}{R_{icr}+W/2}\right),
\end{aligned}
\label{eq:agv_steering_geometry}
\end{equation}
where $L_w$ and $W$ are the wheelbase and track width. Both targets are set to zero when $|\omega_{cmd}|<\omega_{min}$ or $v<v_{low}$. A first-order actuator drives each $\delta_i$ toward $\delta_i^{tgt}$ with time constant $\tau_\delta$, subject to the rate bound $|\dot\delta_i|\leq\dot\delta_{max}$ and angle bound $|\delta_i|\leq\delta_{max}$.

\subsubsection{Dynamic Model}

The motor--gearbox chain limits the drive command before it reaches the wheels. For current limit $I_{max}$, torque constant $k_t$, gear ratio $n$, transmission efficiency $\eta$, and wheel radius $r$, the single-wheel force limit is $F_w^{max}=I_{max}k_t\eta n/\max(r,\varepsilon)$. The equivalent first-order current model uses $\tau_e=L_m/R_m$ and the effective current limit $I_{eff}=\min(I_{max},I_{max}^{volt})$, where $I_{max}^{volt}=\max[0,(V_{bus}-k_e\omega_m)/\max(R_m,\varepsilon)]$ and $\omega_m=|v|n/r$. The command allocation and current-to-force dynamics are
\begin{equation}
\begin{aligned}
F_{cmd}^{eff}&=\operatorname{sat}\!\left(F_{cmd},-2F_w^{max},2F_w^{max}\right),\\
F_{x,lf}^{cmd}&=F_{x,rr}^{cmd}=\frac{F_{cmd}^{eff}}{2},\\
I_i^{tgt}&=\operatorname{sat}\!\left(\frac{F_{x,i}^{cmd}r}{k_tn\eta},-I_{eff},I_{eff}\right),\\
\dot I_i&=\frac{I_i^{tgt}-I_i}{\tau_e},\qquad
F_{x,i}^{pre}=I_i\frac{k_tn\eta}{r},\quad i\in\{lf,rr\}.
\end{aligned}
\label{eq:agv_drive_actuation}
\end{equation}
The equal allocation introduces no differential drive force. The vector $\bm x$ in \eqref{eq:agv_state_control} contains only vehicle motion states. To distinguish it from the state advanced by the simulation, define the full nonlinear plant state as
\begin{equation}
\bm\zeta=
\begin{bmatrix}
\bm x^{\mathsf T}&I_{lf}&I_{rr}&\bm\zeta_f^{\mathsf T}
\end{bmatrix}^{\mathsf T},
\label{eq:full_nonlinear_state}
\end{equation}
where $\bm\zeta_f$ collects the internal states of the implemented actuator and response filters. The currents $I_{lf}$ and $I_{rr}$ are therefore first-order actuator states of the full nonlinear plant rather than elements of the vehicle motion-state vector. Their responses are retained because they provide grade-sensitive measurements for the temporal estimator. The controller introduced later uses a separate four-dimensional path-error state and does not imply that the internal plant states disappear.

The road grade enters the dynamic model through the tangential and normal components of gravity. The grade, rolling, and aerodynamic resistance terms are
\begin{equation}
\begin{aligned}
F_{slope}&=mg\sin\theta,\\
F_{roll}^{est}&=c_rmg\cos\theta,\\
F_{aero}&=\frac{1}{2}\rho_{air}C_DA\,v|v|,\\
F_{res}^{est}&=F_{roll}^{est}+F_{aero}+F_{slope}.
\end{aligned}
\label{eq:agv_slope_resistance}
\end{equation}
where $\rho_{air}$ is the air density and $C_DA$ is the product of the drag coefficient and frontal area. Thus, $mg\sin\theta$ directly changes the longitudinal force balance, whereas $mg\cos\theta$ changes the normal-load-dependent resistance and tire-force capacity. Based on the pre-clipping drive forces, the acceleration estimate, load transfers, and front--rear base loads are
\begin{equation}
\begin{aligned}
a_{long}&=\operatorname{sat}\!\left(
\frac{F_{x,lf}^{pre}+F_{x,rr}^{pre}-F_{res}^{est}}{m},-a_{max},a_{max}\right),\\
\Delta_{long}&=\frac{ma_{long}h_{cg}}{L_w},\qquad
\Delta_{lat}=\frac{m|v\omega|h_{cg}}{\max(W,\varepsilon)},\\
N_f^0&=\frac{mg\cos\theta}{2}-\Delta_{long},\qquad
N_r^0=\frac{mg\cos\theta}{2}+\Delta_{long}.
\end{aligned}
\label{eq:agv_load_transfer}
\end{equation}
Here, $N_f^0$ and $N_r^0$ denote total front- and rear-axle loads, respectively. With the lateral transfer shared equally between the two axles, the individual wheel loads before the nonnegative numerical safeguard are
\begin{equation}
\begin{aligned}
\widetilde N_{lf}&=\frac{N_f^0}{2}-\frac{\operatorname{sgn}(\omega)\Delta_{lat}}{2},&
\widetilde N_{rf}&=\frac{N_f^0}{2}+\frac{\operatorname{sgn}(\omega)\Delta_{lat}}{2},\\
\widetilde N_{lr}&=\frac{N_r^0}{2}-\frac{\operatorname{sgn}(\omega)\Delta_{lat}}{2},&
\widetilde N_{rr}&=\frac{N_r^0}{2}+\frac{\operatorname{sgn}(\omega)\Delta_{lat}}{2}.
\end{aligned}
\label{eq:agv_wheel_load_distribution}
\end{equation}
The implementation projects $\widetilde{\bm N}$ onto the nonnegative simplex with total load $mg\cos\theta$ before calculating $F_{roll}=c_r\sum_iN_i$. Because body-pitch dynamics are excluded, the model neglects the additional steady front--rear load redistribution that a pitched body could produce on a constant grade; $\Delta_{long}$ represents dynamic longitudinal transfer only.

The wheel loads constrain both lateral and longitudinal tire forces. For the two active wheels,
\begin{equation}
\begin{aligned}
\alpha_{lf}&=\delta_{lf}-\left(\beta+\frac{L_w\omega}{2\max(|v|,v_{low})}\right),\\
\alpha_{rr}&=\delta_{rr}-\left(\beta-\frac{L_w\omega}{2\max(|v|,v_{low})}\right),\\
F_{y,i}&=\operatorname{sat}\!\left(-C_{\alpha,i}\alpha_i,-\mu N_i,\mu N_i\right),\\
F_{x,i}^{allow}&=\mu N_i\sqrt{\max\!\left[1-\left(\frac{F_{y,i}}{\max(\mu N_i,\varepsilon)}\right)^2,0\right]},\\
F_{x,i}&=\operatorname{sgn}(F_{x,i}^{pre})\min\!\left(|F_{x,i}^{pre}|,F_{x,i}^{allow}\right),
\quad i\in\{lf,rr\}.
\end{aligned}
\label{eq:agv_tire_forces}
\end{equation}
For linearization, the hard saturation may be replaced by a smooth approximation. Let $\bm{R}(\delta_i)$ denote the planar rotation from wheel to body coordinates and let $\bm{r}_{lf}=[L_w/2,W/2]^T$ and $\bm{r}_{rr}=[-L_w/2,-W/2]^T$ denote the active-wheel positions. The body-frame drive resultants are
\begin{equation}
\begin{aligned}
\bm{F}^{drv}&=\sum_{i\in\{lf,rr\}}\bm{R}(\delta_i)
\begin{bmatrix}F_{x,i}\\0\end{bmatrix},\\
M_z^{drv}&=\sum_{i\in\{lf,rr\}}
\left(\bm{r}_i\times\bm{R}(\delta_i)
\begin{bmatrix}F_{x,i}\\0\end{bmatrix}\right)_z.
\end{aligned}
\label{eq:agv_drive_resultants}
\end{equation}

Finally, let $m_{eff}=m+2(I_w+I_mn^2)/r^2$ account for the reflected wheel and motor inertia. Using the resistance terms in \eqref{eq:agv_slope_resistance}, the longitudinal, yaw, and sideslip dynamics are
\begin{equation}
\begin{aligned}
\dot v&=\frac{F_x^{drv}-F_{roll}-F_{aero}-F_{slope}}
{\max(m_{eff},\varepsilon)},\\
M_z&=\frac{L_w}{2}(F_{y,lf}-F_{y,rr})+M_z^{drv}-C_{yaw}\omega,\\
\dot\omega&=\operatorname{sat}\!\left(
\frac{M_z}{\max(I_z,\varepsilon)},-\dot\omega_{max},\dot\omega_{max}\right),\\
\dot\beta&=\frac{F_{y,lf}+F_{y,rr}+F_y^{drv}}
{m_{eff}\max(|v|,v_{low}/10)}-\omega-C_{\beta}\beta.
\end{aligned}
\label{eq:agv_continuous_dynamics}
\end{equation}
At low speed, the sideslip equation is replaced by $\dot\beta=-C_{\beta}^{low}\beta$. The yaw-acceleration and sideslip-rate limits are $\dot\omega_{max}=5~\mathrm{rad/s^2}$ and $|\dot\beta|\leq10^{\circ}/\mathrm{s}$, respectively. Thus, $\theta$ does not directly change the planar kinematic mapping in \eqref{eq:agv_kinematics}; instead, $mg\sin\theta$ changes the longitudinal force balance and the $\cos\theta$-dependent wheel loads modify the available tire forces. These effects propagate through $v$, $\omega$, and $\beta$ to the path coordinates and provide the grade-sensitive onboard responses used by the temporal estimator. The parameter values that define the plant are listed in Table~\ref{tab:agv_params}; the same values are shared by data generation, LPV linearization, and closed-loop evaluation.

\begin{table}[!t]
\centering
\caption{\textbf{Physical parameters of the diagonal dual-steering-wheel AGV plant.}}
\label{tab:agv_params}
\setlength{\tabcolsep}{4.5pt}
\setlength{\tabcolsep}{3pt}
\begin{tabular}{|l|l|r|}
\hline
Symbol & Description & Value \\
\hline
$m$ & Vehicle mass & $200~\mathrm{kg}$ \\
$I_z$ & Yaw moment of inertia & $16.67~\mathrm{kg\,m^2}$ \\
$L_w$ & Wheelbase & $2.0~\mathrm{m}$ \\
$W$ & Track width & $0.8~\mathrm{m}$ \\
$h_{cg}$ & Center-of-gravity height & $0.5~\mathrm{m}$ \\
\hline
$r$ & Wheel radius & $0.15~\mathrm{m}$ \\
$I_w$ & Wheel rotational inertia & $0.0135~\mathrm{kg\,m^2}$ \\
$I_m$ & Motor-rotor inertia & $1.77\times10^{-3}~\mathrm{kg\,m^2}$ \\
$k_t$ & Motor torque constant & $1.21~\mathrm{N\,m/A}$ \\
$n$ & Gear ratio & $10$ \\
$\eta$ & Transmission efficiency & $0.9$ \\
$I_{max}$ & Per-motor current limit & $9.0~\mathrm{A}$ \\
\hline
$\tau_\delta$ & Steering time constant & $0.08~\mathrm{s}$ \\
$\delta_{max}$ & Steering-angle limit & $90^{\circ}$ \\
$\dot\delta_{max}$ & Steering-rate limit & $30^{\circ}/\mathrm{s}$ \\
$a_{max}$ & Longitudinal-acceleration bound & $1.0~\mathrm{m/s^2}$ \\
\hline
$C_{\alpha}$ & Cornering stiffness (each active wheel) & $3000~\mathrm{N/rad}$ \\
$\mu$ & Tire--road friction coefficient & $0.8$ \\
$c_r$ & Rolling-resistance coefficient & $0.015$ \\
$\rho_{air}$ & Air density & $1.225~\mathrm{kg/m^3}$ \\
$C_DA$ & Drag-area product & $0.5~\mathrm{m^2}$ \\
$C_{yaw}$ & Yaw damping coefficient & $250~\mathrm{N\,m\,s/rad}$ \\
$C_{\beta}$, $C_{\beta}^{low}$ & Sideslip damping (normal, low speed) & $0$, $1.0~\mathrm{s^{-1}}$ \\
\hline
\end{tabular}
\end{table}


\subsection{Slope-Aware Path-Tracking Objective and Problem Formulation}

The objective of this work is constrained path tracking under longitudinal
grade variations, rather than grade estimation as an end in itself; grade coupling has likewise motivated dedicated adaptive tracking designs for multi-wheel-drive platforms \citep{Zhang2025CoupledSlope4WID}. The
tracking error is defined as
$\bm e_k=[e_{y,k},e_{\psi,k},e_{v,k},e_{\omega,k}^{\mathrm{trk}}]^{T}$, which
collects the lateral, heading, longitudinal-speed, and yaw-rate errors. The
superscript distinguishes $e_{\omega}^{\mathrm{trk}}$ from the kinematic
yaw-consistency feature introduced in
Section~\ref{sec:proposed_framework}. The control problem is stated in the
general form
\begin{equation}
\begin{aligned}
\min_{\{\bm u_k\}}\quad &J_{\mathrm{track}}(\bm e,\Delta\bm u)\\
\mathrm{s.t.}\quad
&\bm\zeta_{k+1}=\mathcal F_{T_s}(\bm\zeta_k,\bm u_k,\theta_k),\\
&\bm u_k\in\mathcal U,\qquad
\Delta\bm u_k\in\Delta\mathcal U,
\end{aligned}
\label{eq:slope_aware_tracking_objective}
\end{equation}
where $\mathcal F_{T_s}(\cdot)$ is the one-step transition of the AGV model,
and $\mathcal U$ and $\Delta\mathcal U$ are the admissible input and
input-increment sets. The finite-horizon cost and constraints are detailed
with the controller formulation in Section~\ref{sec:model_lpv_mpc}. Within
this problem, the grade angle is a control-relevant variable rather than an
isolated geometric quantity. It changes the longitudinal force balance
through $mg\sin\theta_k$ in \eqref{eq:agv_slope_resistance} and modifies the
available tire forces through the $\cos\theta_k$-dependent wheel loads in
\eqref{eq:agv_tire_forces}. Scheduling the controller with an inaccurate
grade therefore makes the prediction model inconsistent with the plant and
degrades closed-loop tracking
performance~\citep{Sato2013InexactScheduling,Metzler2021PredictionFidelity}.

A direct route to $\theta_k$ is inertial sensing. However, IMU-derived grade
estimates can be biased by transient longitudinal acceleration, body-pitch
excitation, road-induced vibration, and low-cost sensor
noise~\citep{Liu2023RoadSlopeAAIMM}. Model-based observers that fuse
inertial and drivetrain information also remain sensitive to
vehicle-parameter mismatch and load
variation~\citep{Gao2024RoadSlopeHDV}. The controller is therefore assumed
to have access to neither the true grade nor a reliable instantaneous
inertial measurement. Instead, the current grade is inferred from a fixed
window of onboard responses available at the sampling instant,
\begin{equation}
\mathcal H_k=\left\{\bm s_{k-N_h+1},\bm s_{k-N_h+2},\ldots,\bm s_k\right\},
\qquad
\hat\theta_k^{\mathrm{hist}}=f_{\phi}(\mathcal H_k),
\label{eq:slope_estimation_objective}
\end{equation}
where $\bm s_k$ collects measured or recursively derived proprioceptive
quantities. The true grade and any direct pitch-based grade measurement are
excluded from the estimator input, and no sample later than the current
instant $k$ is used. The window length $N_h$ balances the observability of
grade-induced responses against the dilution of grade-transition
information. Its numerical value and the underlying response-time
calculation are reported with the experimental settings.

These considerations reduce slope-aware tracking to four coupled
requirements: a lightweight history-based temporal estimator,
dynamics-informed feature and multi-scale lag construction,
uncertainty-adaptive fusion of the history-based estimate and the qualified observer estimate, and
synchronized incorporation of the fused grade into all grade-dependent
LPV-MPC updates. Section~\ref{sec:proposed_framework} develops these
components in turn to support the tracking objective in
\eqref{eq:slope_aware_tracking_objective}.

\begin{figure*}[!t]
\centering
\includegraphics[width=\textwidth]{../7.6/figures_final/Fig01_overall_framework/output/fig01_overall_framework.pdf}
\caption{Overall slope-aware perception-and-control framework. Arrows trace the online signal flow from the onboard measurements through the temporal-perception, uncertainty-adaptive fusion, and grade-scheduled LPV-MPC stages to the AGV plant; band headers indicate the sections specifying each stage. The dashed box marks the offline stage, whose dataset, trained weights, and fusion calibration are deployed before online operation.}
\label{fig:overall_framework}
\end{figure*}

\section{Slope-Aware Temporal Perception and LPV-MPC Framework}
\label{sec:proposed_framework}

The formulation in Section~\ref{sec:agv_slope_estimation} creates a perception-to-control chain rather than four independent additions. The first problem is to infer a control-relevant grade from onboard histories without using a direct grade channel. The second is to expose both the absolute operating state and the short grade-sensitive changes that an instantaneous input can obscure. The third is to reconcile the complementary but condition-dependent information supplied by the temporal estimator and the qualified observer. The final problem is to inject one physically consistent grade signal into every grade-dependent LPV-MPC update. Accordingly, the proposed framework proceeds from a lightweight ModernTCN estimator, through dynamics-informed feature and delta-lag construction, to adaptive grade fusion and synchronized LPV-MPC scheduling. Figure~\ref{fig:overall_framework} summarizes this pipeline and its interfaces: onboard measurement histories feed the temporal-perception branch and the qualified observer branch in parallel, a single fused grade $\theta_k^{\mathrm{sch}}$ crosses the perception-to-control boundary, and every grade-dependent element of the LPV-MPC update is scheduled from this signal synchronously.

\subsection{Lightweight ModernTCN-Based Grade Estimator}

To implement the history-to-grade mapping in \eqref{eq:slope_estimation_objective}, a lightweight task-specific estimator is developed from the temporal modeling principles of ModernTCN~\citep{Luo2024ModernTCN}. The estimator retains large-kernel depthwise temporal convolution, pointwise channel mixing, and residual learning, while removing the full patch-based and multi-stage hierarchy of the original general-purpose architecture. This reduction gives a fixed-window network that can be exported through ONNX and executed within the MATLAB/Simulink closed loop.

Let the network input be written generically as
\begin{equation}
\mathbf Z_k\in\mathbb R^{N_h\times F},
\label{eq:generic_moderntcn_input}
\end{equation}
where its physical channels and temporal augmentation are specified in the following two subsections. A pointwise stem first projects the input to a latent channel space,
\begin{equation}
\mathbf H_k^{(0)}=\sigma\!\left[
\operatorname{BN}\!\left(
\operatorname{Conv}_{1\times1}(\mathbf Z_k^{\mathsf T})
\right)\right],
\label{eq:moderntcn_stem}
\end{equation}
where $\mathbf H_k^{(0)}\in\mathbb R^{C\times N_h}$, $C=64$, $\operatorname{BN}(\cdot)$ denotes batch normalization, and $\sigma(\cdot)$ is the ReLU activation. Five stride-one residual temporal blocks are then applied. The $\ell$th block is
\begin{align}
\mathbf U_k^{(\ell)}
 &=\operatorname{DWConv}_{K}\!\left(\mathbf H_k^{(\ell)}\right),\nonumber\\
\mathbf V_k^{(\ell)}
 &=\operatorname{PW}_{2}^{(\ell)}\!\left\{
\sigma\!\left[\operatorname{PW}_{1}^{(\ell)}
\!\left(\sigma\!\left[\operatorname{BN}
(\mathbf U_k^{(\ell)})\right]\right)\right]\right\},\nonumber\\
\mathbf H_k^{(\ell+1)}
 &=\mathbf H_k^{(\ell)}
 +\gamma_{\ell}\operatorname{BN}\!\left(\mathbf V_k^{(\ell)}\right),
\label{eq:moderntcn_small_block}
\end{align}
where $\operatorname{DWConv}_{K}$ is a depthwise temporal convolution with $K=31$, the two pointwise convolutions use an expansion ratio of two, and $\gamma_{\ell}$ is a learnable residual scale. The five large-kernel blocks provide a nominal receptive field of 151 samples, thereby covering the implemented historical window.

The final sequence is summarized by its last, mean, and maximum latent responses. Window-level statistics of the normalized input are added to retain low-order information about each channel:
\begin{equation}
\begin{aligned}
\mathbf h_k={}&\left[
\mathbf H_{k,:,N_h}^{(5)};
\operatorname{mean}_{t}\mathbf H_{k,:,t}^{(5)};
\operatorname{max}_{t}\mathbf H_{k,:,t}^{(5)};
q(\mathbf Z_k)\right],\\
\hat{\theta}^{\mathrm{raw}}_k
 ={}&\mathbf w_{\theta}^{\mathsf T}\mathbf h_k+b_{\theta},
\end{aligned}
\label{eq:moderntcn_readout}
\end{equation}
where $q(\mathbf Z_k)$ concatenates the current value, mean, standard deviation, maximum, and minimum of each input channel. Although same padding is used inside the fixed window, $\mathbf Z_k$ terminates at the current sample and contains no measurement later than $k$. The deployed estimator is therefore a history-only online mapping; it is not described as a strictly causal convolutional backbone. Before fusion, the raw regression output is processed by a causal deployment interface, which applies first-order temporal smoothing, the configured physical magnitude and rate limits, and the MPC-side deadband. The resulting control-facing estimate is denoted by $\hat{\theta}^{\mathrm{MTCN}}_k$ in the remainder of this section. This distinction prevents the network readout from being conflated with the stateful signal actually consumed by the fusion module.

The shared representation additionally feeds two auxiliary classification heads,
\begin{equation}
\begin{aligned}
\hat{\mathbf p}^{\mathrm{main}}_k
&=\operatorname{softmax}\!\left(g_{\mathrm{main}}(\mathbf h_k)\right),\\
\hat{\mathbf p}^{\mathrm{turn}}_k
&=\operatorname{softmax}\!\left(g_{\mathrm{turn}}(\mathbf h_k)\right),
\end{aligned}
\label{eq:slope_auxiliary_outputs}
\end{equation}
which identify the main operating condition and steering direction. Both heads are auxiliary to grade regression and are not evaluated as independent perception contributions. The steering-direction output is retained for representation learning and diagnostics; only $c_k^{\mathrm{main}}=\max(\hat{\mathbf p}^{\mathrm{main}}_k)$ enters the deployed interface, where it calibrates the ModernTCN uncertainty used by the fusion weight. Neither predicted label directly schedules the controller or acts as a hard switching predicate. The training objective is
\begin{equation}
\mathcal L=
\lambda_{\theta}\mathcal L_{\theta}
+\mathcal L_{\mathrm{main}}
+\lambda_{\mathrm{turn}}\mathcal L_{\mathrm{turn}}
+\mathcal R_{\theta},
\label{eq:slope_estimator_loss}
\end{equation}
where $\mathcal R_{\theta}$ collects the flat-road and bounded-error grade penalties. Table~\ref{tab:moderntcn_variant_comparison} distinguishes the implemented network from the complete ModernTCN architecture.

\begin{table*}[!t]
\centering
\caption{\textbf{Structural distinction between the original ModernTCN and the lightweight task-specific estimator.}}
\label{tab:moderntcn_variant_comparison}
\setlength{\tabcolsep}{3pt}
\begin{tabular}{|p{0.20\textwidth}|p{0.35\textwidth}|p{0.37\textwidth}|}
\hline
Component & Original ModernTCN~\citep{Luo2024ModernTCN} & Lightweight task-specific estimator used in this work \\
\hline
Input embedding & Variable-independent patch embedding with an explicit variable dimension & Joint $1\times1$ channel projection of the fixed historical input \\
Temporal modeling & Large- and small-kernel depthwise convolutions with structural re-parameterization & One large-kernel depthwise convolution ($K=31$) in each residual block \\
Feature mixing & Decoupled grouped ConvFFN modules for within- and cross-variable interaction & Two standard pointwise convolutions with an expansion ratio of two \\
Backbone hierarchy & Patch tokens and a multi-stage representation hierarchy & Five repeated stride-one residual blocks with 64 latent channels \\
Readout and deployment & Task-dependent pipelines for general time-series analysis & Last/mean/max latent pooling plus input statistics, grade and auxiliary heads, and ONNX deployment in MATLAB/Simulink \\
\hline
\end{tabular}
\end{table*}

\subsection{Dynamics-Informed 22-D Proprioceptive Feature Construction}

The temporal architecture alone does not determine whether the input contains identifiable grade information. The AGV model shows that grade first changes longitudinal resistance and available tire forces, after which its effects appear jointly in speed, acceleration, drive effort, wheel asymmetry, and yaw--steering consistency. The estimator input is therefore restricted to online onboard measurements and recursively computed response variables that represent these complementary mechanisms, consistent with longitudinal-dynamics-based estimation from standard IMU and CAN-bus signals \citep{Jensen2022MassGradeIMUCAN}. The resulting fixed feature vector is
\begin{equation}
\begin{aligned}
\mathbf{s}_k=\big[&
\omega_{g},I_{lf},I_{rr},\omega_{w,lf},\omega_{w,rr},
\delta_{lf},\delta_{rr},\\
&\hat v,\dot{\hat v},\Delta\omega_w,I_{\Sigma},
\Delta I,\Delta I_{|\cdot|},\\
&\hat\kappa,a_I,\dot{\hat v}_{\mathrm{LP}},a_{x,w},
I_{\mathrm{drv}},I_a,\Lambda,a_{\mathrm{HP}},\varepsilon_{\omega}^{\mathrm{kin}}
\big]_k^{\mathsf T}\in\mathbb R^{22}.
\end{aligned}
\label{eq:feature_vector_22d}
\end{equation}

Here, $\omega_g$ is the measured yaw rate; $I_{lf}$ and $I_{rr}$ are the active-wheel motor currents; $\omega_{w,lf}$ and $\omega_{w,rr}$ are the corresponding wheel speeds; and $\delta_{lf}$ and $\delta_{rr}$ are the steering angles. The wheel-speed-derived motion responses are updated recursively as
\begin{align}
\hat v_k&=\frac{r}{2}\left(\omega_{w,lf,k}+\omega_{w,rr,k}\right),\nonumber\\
a_{x,w,k}&=\frac{\hat v_k-\hat v_{k-1}}{T_s},\nonumber\\
\dot{\hat v}_k&=\alpha_d a_{x,w,k}
+(1-\alpha_d)\dot{\hat v}_{k-1},\nonumber\\
\dot{\hat v}_{\mathrm{LP},k}
&=\alpha_{\mathrm{LP}}\dot{\hat v}_k
+(1-\alpha_{\mathrm{LP}})\dot{\hat v}_{\mathrm{LP},k-1},\nonumber\\
a_{\mathrm{HP},k}&=a_{x,w,k}-\dot{\hat v}_{\mathrm{LP},k},
\label{eq:kinematic_feature_channels}
\end{align}
where $\alpha_d=T_s/(\tau_d+T_s)$ and $\alpha_{\mathrm{LP}}=T_s/(\tau_{\mathrm{LP}}+T_s)$. The raw, filtered, and high-pass terms retain rapid wheel-speed changes, the persistent acceleration trend, and short longitudinal transients, respectively.

At the low operating speeds considered here, $\hat v$ is used only as a wheel-speed-derived motion proxy. It does not correct for steering projection or longitudinal slip and is therefore not interpreted as an exact vehicle speed during sharp steering or tire slip. The bilateral drive responses are
\begin{align}
\Delta\omega_{w,k}&=\left|\omega_{w,lf,k}-\omega_{w,rr,k}\right|,\nonumber\\
I_{\Sigma,k}&=|I_{lf,k}|+|I_{rr,k}|,\qquad
\Delta I_k=I_{lf,k}-I_{rr,k},\nonumber\\
\Delta I_{|\cdot|,k}&=|I_{lf,k}|-|I_{rr,k}|,\qquad
I_{\mathrm{drv},k}=I_{lf,k}+I_{rr,k},
\label{eq:drive_asymmetry_channels}
\end{align}
which distinguish total traction demand from unequal wheel loading and turning-related allocation. The remaining consistency variables are
\begin{align}
\hat\kappa_k&=\frac{\tan\delta_{lf,k}-\tan\delta_{rr,k}}{L_w},\nonumber\\
a_{I,k}&=\frac{\dot{\hat v}_k}{\max(I_{\Sigma,k},I_{\min})},\qquad
I_{a,k}=\frac{I_{\Sigma,k}}
{\max(|\dot{\hat v}_{\mathrm{LP},k}|,a_{\min})},\nonumber\\
\Lambda_k&=
\frac{I_{\mathrm{drv},k}-\mu_{I,\mathrm{tr}}}{\max(\sigma_{I,\mathrm{tr}},\epsilon_N)}
-\frac{\dot{\hat v}_{\mathrm{LP},k}-\mu_{a,\mathrm{tr}}}{\max(\sigma_{a,\mathrm{tr}},\epsilon_N)},\nonumber\\
\varepsilon_{\omega,k}^{\mathrm{kin}}&=\omega_{g,k}-\hat v_k\hat\kappa_k .
\label{eq:load_consistency_channels}
\end{align}
The protected denominators $I_{\min}=0.1$ and $a_{\min}=0.05$ prevent singular behavior at low current or acceleration. The statistics $(\mu_{I,\mathrm{tr}},\sigma_{I,\mathrm{tr}})$ and $(\mu_{a,\mathrm{tr}},\sigma_{a,\mathrm{tr}})$ are fitted on the training split, and $\epsilon_N>0$ protects the standardization. Thus, $\Lambda_k$ combines two separately standardized quantities and is dimensionless; it is interpreted as an empirical drive--response co-variation feature rather than as a physical force. The denominator $L_w$ in $\hat\kappa_k$ follows the front--rear steering geometry in \eqref{eq:agv_steering_geometry}; the quantity remains a low-speed curvature proxy because slip and passive-wheel kinematics are omitted. Finally, $\varepsilon_{\omega,k}^{\mathrm{kin}}$ expresses yaw inconsistency and is distinct from the MPC tracking error $e_{\omega,k}^{\mathrm{trk}}$. The complete 22-channel input contract assembled from these mechanisms is listed in Table~\ref{tab:feature_contract_22d}.

\begin{table*}[!t]
\centering
\caption{\textbf{Dynamics-informed 22-dimensional proprioceptive feature set.}}
\label{tab:feature_contract_22d}
\setlength{\tabcolsep}{3pt}
\begin{tabular}{|>{\raggedright\arraybackslash}p{0.16\textwidth}|>{\centering\arraybackslash}p{0.05\textwidth}|>{\raggedright\arraybackslash}p{0.39\textwidth}|>{\raggedright\arraybackslash}p{0.28\textwidth}|}
\hline
Feature group & Dim. & Implemented channels & Primary information represented \\
\hline
Raw onboard sensing & 7 & \texttt{gyro\_z}, \texttt{I\_lf}, \texttt{I\_rr}, \texttt{omega\_wheel\_lf}, \texttt{omega\_wheel\_rr}, \texttt{delta\_lf}, \texttt{delta\_rr} & Measured yaw, drive effort, wheel motion, and steering state \\
Kinematic and multi-band response & 5 & \texttt{v\_hat}, \texttt{dv\_hat\_dt}, \texttt{dv\_hat\_dt\_lp}, \texttt{accel\_x\_wheel}, \texttt{a\_hp} & Current speed, persistent acceleration trend, and rapid longitudinal transients \\
Drive asymmetry and allocation & 5 & \texttt{ws\_imbalance}, \texttt{I\_sum}, \texttt{I\_diff\_signed}, \texttt{I\_diff\_abs}, \texttt{I\_drive\_signed} & Total traction demand and unequal active-wheel response \\
Load and kinematic consistency proxies & 5 & \texttt{kappa\_proxy}, \texttt{accel\_per\_current}, \texttt{current\_per\_accel}, \texttt{drive\_load\_proxy}, \texttt{yaw\_consistency\_error} & Steering curvature, drive--acceleration coupling, and yaw inconsistency \\
\hline
\end{tabular}
\end{table*}

The feature set excludes global pose, unmeasured sideslip and wheel normal loads, the true road grade, and direct pitch- or grade-measurement channels. The yaw-rate gyroscope remains because it measures rotational vehicle response rather than grade. Thus, the network must infer grade from the coupled response pattern of the AGV instead of learning an identity mapping from a direct grade measurement.

\subsection{Multi-Scale Delta-Lag Representation and Historical Window Assembly}

Even with the 22-D feature set, a single sample is ambiguous because similar current and acceleration levels can be produced by grade resistance, a traction or braking command, or a steering transition. Sequence models over onboard response histories have proven effective for related road-condition identification tasks \citep{Liang2022LSTMRoadID}. Temporal differences are therefore constructed after train-set normalization,
\begin{equation}
\mathbf z_k=\mathcal N(\mathbf s_k)
=\frac{\mathbf s_k-\boldsymbol\mu_{\mathrm{tr}}}
{\boldsymbol\sigma_{\mathrm{tr}}},
\label{eq:feature_train_normalization}
\end{equation}
where $\boldsymbol\mu_{\mathrm{tr}}$ and $\boldsymbol\sigma_{\mathrm{tr}}$ are fitted only on the training split. For a lag $\ell$, the normalized delta channel is
\begin{equation}
\Delta_{\ell}\mathbf z_k=
\operatorname{clip}\!\left(
\mathbf z_k-\mathbf z_{k-\ell},-c_{\Delta},c_{\Delta}
\right),\qquad \ell\in\{1,2,4\},
\label{eq:multiscale_delta_lag}
\end{equation}
with $c_{\Delta}=5$. At $T_s=0.01$ s, the three offsets represent 10, 20, and 40 ms response changes. The augmented vector is
\begin{equation}
\widetilde{\mathbf z}_k=
\begin{bmatrix}
\mathbf z_k^{\mathsf T}&
(\Delta_1\mathbf z_k)^{\mathsf T}&
(\Delta_2\mathbf z_k)^{\mathsf T}&
(\Delta_4\mathbf z_k)^{\mathsf T}
\end{bmatrix}^{\mathsf T}
\in\mathbb R^{88}.
\label{eq:delta_bank_124}
\end{equation}

The augmented representation in \eqref{eq:delta_bank_124} is referred to as the delta bank. The current normalized state retains the absolute operating condition, whereas the three difference blocks expose increasingly wider local changes. Calculating the differences after normalization prevents the wheel-speed and current scales from dominating the augmentation, and clipping limits isolated transients. The lag set was configured through a staged representation study: single offsets from $\{1,2,4,8,16,32,64\}$ were first screened, after which the short offsets $\{1,2,4\}$ were carried forward. Raw input, a single-difference branch, multi-lag state stacking, the three-branch delta bank, and mixed lag--delta inputs were then compared under the same estimator and synchronized controller interface. The delta bank in \eqref{eq:delta_bank_124} is the deployed representation selected from this matched closed-loop ablation; the multi-seed comparison is reported with the experimental results.

The 88-dimensional vectors are placed in chronological order,
\begin{equation}
\mathbf Z_k=
\begin{bmatrix}
\widetilde{\mathbf z}_{k-N_h+1}^{\mathsf T}\\
\widetilde{\mathbf z}_{k-N_h+2}^{\mathsf T}\\
\vdots\\
\widetilde{\mathbf z}_{k}^{\mathsf T}
\end{bmatrix}
\in\mathbb R^{N_h\times88},\qquad N_h=128.
\label{eq:historical_input_128x88}
\end{equation}
At the 0.01-s sampling interval, the implemented window covers 1.28 s. Its length is fixed from the characteristic response time over which grade-induced resistance produces an observable change in the onboard responses at the representative speed, grade, and effective inertia of the present AGV; the corresponding parameter calculation is given in the experimental setup. For unavailable lagged samples at the beginning of a window, the first normalized sample is used for edge padding, which makes the corresponding differences zero. The same normalization, padding, endpoint labeling, and channel order are used during dataset construction, validation, ONNX export, and online inference.

\subsection{Uncertainty-Adaptive Grade Fusion with the Qualified Observer}
\label{sec:adaptive_slope_fusion}

The history-based estimator and the qualified observer provide complementary grade information. Combining heterogeneous slope estimates has previously been addressed through multi-model and multi-source fusion \citep{Feng2021MultiModelFusion}; here, the combination must additionally respect the condition-dependent reliability of a learned source. ModernTCN infers grade from coupled vehicle histories without a direct pitch or grade channel, whereas the observer uses inertial specific force and angular rate. The proposed interface therefore retains ModernTCN as the primary source and admits only an uncertainty-weighted, bounded correction from the qualified observer.

The observer follows the robust Kalman filter with external-acceleration compensation in \citep{Chen2024RobustAttitude}; its full state-space derivation is not repeated here. The deployed input is the bias-compensated six-axis IMU packet $[\tilde{\bm a}_k^{T},\tilde{\bm\omega}_k^{T}]^{T}$. Its unit gravity-direction state $\bm x_k=[C_{31,k},C_{32,k},C_{33,k}]^{T}$ and accelerometer innovation are propagated as
\begin{equation}
\begin{aligned}
\bm\Phi_{k-1}&=\bm I_3+T_s[\tilde{\bm\omega}_k]_{\times},\\
\bm x_k^{-}&=\operatorname{unit}(\bm\Phi_{k-1}\bm x_{k-1}),\\
\bm P_k^{-}&=\bm\Phi_{k-1}\bm P_{k-1}\bm\Phi_{k-1}^{T}+\bm Q_{k-1},\\
\hat{\bm a}_{k}^{e,-}&=c_a\hat{\bm a}_{k-1}^{e},\qquad \bm H_k=g\bm I_3,\\
\bm r_k^{a}&=\tilde{\bm a}_k-\hat{\bm a}_{k}^{e,-}-\bm H_k\bm x_k^{-},
\end{aligned}
\label{eq:qualified_observer_prediction}
\end{equation}
where $\operatorname{unit}(\bm z)=\bm z/\|\bm z\|_2$, $[\cdot]_{\times}$ is the cross-product matrix, $\bm Q_{k-1}$ is the gyroscope-noise-induced process covariance, and $c_a=\exp(-T_s/0.25)$ is the external-acceleration persistence coefficient. From a window of $\mu$ innovations, the observer estimates the external-acceleration covariance and performs the adaptive correction,
\begin{equation}
\begin{aligned}
\hat{\bm C}_{r,k}&=\frac{1}{\mu}\sum_{j=k-\mu+1}^{k}
\bm r_j^{a}(\bm r_j^{a})^{T},\\
\bm D_k&=\hat{\bm C}_{r,k}-\bm H_k\bm P_k^{-}\bm H_k^{T}-\bm\Sigma_a,\\
\hat{\bm\Sigma}_{k}^{e,*}&=\operatorname{diag}\!\left(
\operatorname{clip}\!\left[
\operatorname{diag}(\bm D_k),
0,\sigma_{e,\max}^{2}\right]\right),\\
\bm R_k&=\bm\Sigma_a+\hat{\bm\Sigma}_{k}^{e},\\
\bm K_k&=\bm P_k^{-}\bm H_k^{T}
\left(\bm H_k\bm P_k^{-}\bm H_k^{T}+\bm R_k\right)^{-1},\\
\bm x_k&=\operatorname{unit}\!\left(\bm x_k^{-}+\bm K_k\bm r_k^{a}\right),\\
\hat{\bm a}_{k}^{e}=\tilde{\bm a}_k-g\bm x_k.
\end{aligned}
\label{eq:qualified_observer_adaptation}
\end{equation}
Here $\bm\Sigma_a$ is the nominal accelerometer covariance and $\sigma_{e,\max}^{2}$ caps the adaptive variance. The candidate $\hat{\bm\Sigma}_{k}^{e,*}$ replaces $\hat{\bm\Sigma}_{k-1}^{e}$ only when the innovation-energy gate in \citep{Chen2024RobustAttitude} detects external acceleration. The covariance state is updated in Joseph form as in the same reference. The resulting road-grade interface is
\begin{equation}
\hat\theta_k^{\mathrm{obs}}
=\operatorname{clip}\!\left[
\operatorname{atan2}(x_{1,k},x_{3,k}),
-\theta_{\max}^{\mathrm{RKF}},\theta_{\max}^{\mathrm{RKF}}
\right].
\label{eq:rkf_imu_grade_output}
\end{equation}
This conversion assumes a calibrated, road-aligned rigid body in low-dynamic motion. The deployed settings are $\mu=10$, $\theta_{\max}^{\mathrm{RKF}}=20^{\circ}$, and a 1-s stationary initialization. The observer is valid only after initialization when all six IMU channels and filter states are finite and $\bm P_k$ is positive semidefinite. It requires no grade truth, wheel speed, drive current, or vehicle force parameter online. The fusion stage consumes only the validity flag, $\hat\theta_k^{\mathrm{obs}}$, and the disturbance statistics defined below.

Because neither deployed source provides a calibrated total prediction-error variance, the fusion statistics are estimated offline from the recorded training and validation runs. Let $e_k^{\mathrm{M}}=\hat\theta_k^{\mathrm{MTCN}}-\theta_k$ and $e_k^{\mathrm{O}}=\hat\theta_k^{\mathrm{obs}}-\theta_k$. Ground-truth grade is used only for this offline calibration and is absent from the deployed interface. Samples are assigned to the same four flat/graded and steady/transition regimes defined by the $1.3^{\circ}$ magnitude and $1^{\circ}\,\mathrm{s}^{-1}$ rate thresholds. Within each regime, the instantaneous disturbance severity of the qualified observer is summarized by
\begin{equation}
q_k^{\mathrm{RKF}}=
\max\!\left\{
\frac{\operatorname{tr}(\hat{\bm\Sigma}_{k}^{e})}
{\operatorname{tr}(\bm\Sigma_a)},
\frac{1}{3}(\bm r_k^{a})^{T}
\left(\bm H_k\bm P_k^{-}\bm H_k^{T}+\bm R_k\right)^{-1}
\bm r_k^{a}
\right\}.
\label{eq:fusion_rkf_disturbance_score}
\end{equation}
Two thresholds frozen from the training-run distribution define low-, medium-, and high-disturbance bins, giving a four-regime by three-disturbance calibration table. Ground-truth grade is not used to compute $q_k^{\mathrm{RKF}}$.

Errors are Winsorized at $\pm5$ robust standard deviations about their per-case median. The four-regime reference moments, estimated from the recorded training and validation data, are the total second moments
\begin{equation}
\begin{aligned}
P_{\mathrm{M},r}^{0}&=\mathbb E[(e^{\mathrm{M}})^2], &
P_{\mathrm{O},r}^{0}&=\mathbb E[(e^{\mathrm{O}})^2],\\
C_{r}^{0}&=\mathbb E[e^{\mathrm{M}}e^{\mathrm{O}}].
\end{aligned}
\label{eq:fusion_empirical_moments}
\end{equation}
This operation limits outliers without removing empirical bias from the total moments. $P_{\mathrm{O},r}^{0}$ is estimated from the qualified-observer outputs recorded on the training runs, whereas $P_{\mathrm{M},r}^{0}$ and $C_r^{0}$ are estimated from the recorded validation runs of the ten trained models. Regime estimates are shrunk toward their global moments and subjected to a fixed variance floor. To introduce condition dependence without replacing these reference moments, the train-only conditional moment of the qualified observer in disturbance cell $(r,q)$ is divided by its sample-exposure-weighted regime mean. The ratio is clipped to $[0.8,1.2]$ and constrained to be nondecreasing from low to high disturbance, giving multiplier $m_{r,q}$. Runtime table entries are then
\begin{equation}
\begin{aligned}
P_{\mathrm{O},r,q}&=m_{r,q}P_{\mathrm{O},r}^{0},\\
C_{r,q}&=\sqrt{m_{r,q}}\,C_r^{0}.
\end{aligned}
\label{eq:fusion_excitation_moments}
\end{equation}
The centered cross-covariance is scaled identically and retained for verification, whereas $C_{r,q}$ is the total cross-error moment used by the online gain. Error correlation is clipped to $[-0.95,0.95]$, and every resulting $2\times2$ total-error matrix is required to be positive semidefinite.

For main-condition confidence $c_k^{\mathrm{main}}$, define $a_k=1+(1-c_k^{\mathrm{main}})$ and
\begin{equation}
P_{\mathrm{M},k}=a_kP_{\mathrm{M},r}^{0},\qquad
P_{\mathrm{O},k}=P_{\mathrm{O},r,q},\qquad
C_k=\sqrt{a_k}\,C_{r,q}.
\label{eq:moderntcn_regime_uncertainty}
\end{equation}
Thus, the auxiliary confidence changes the strength of a possible correction from the qualified observer without acting as a hard switch. With innovation $\nu_k=\hat\theta_k^{\mathrm{obs}}-\hat\theta_k^{\mathrm{MTCN}}$, the innovation variance, normalized innovation squared (NIS), and unconstrained gain are
\begin{equation}
\begin{aligned}
S_k&=\max(P_{\mathrm{M},k}+P_{\mathrm{O},k}-2C_k,S_{\min}),\\
\mathrm{NIS}_k&=\frac{\nu_k^2}{S_k},\\
K_k^{\mathrm{raw}}&=\frac{P_{\mathrm{M},k}-C_k}{S_k}.
\end{aligned}
\label{eq:adaptive_imu_gain}
\end{equation}
where $S_{\min}=(0.02\pi/180)^2$. The innovation weight and effective gain are
\begin{equation}
\begin{aligned}
\eta_k&=
\begin{cases}
1, & \mathrm{NIS}_k\leq9,\\
\sqrt{9/\mathrm{NIS}_k}, &9<\mathrm{NIS}_k\leq25,\\
0, & \mathrm{NIS}_k>25,
\end{cases}\\
K_k^{\mathrm{eff}}&=\eta_k\min[K_{\max},\max(0,K_k^{\mathrm{raw}})].
\end{aligned}
\label{eq:innovation_conditioned_imu_gain}
\end{equation}
with $K_{\max}=0.5$, selected from $\{0.3,0.4,0.5\}$ by a pre-specified validation-based selection rule.

For a valid sample, the target correction and its causal rate-limited update are
\begin{align}
\delta\theta_k^{\star}
&=\operatorname{clip}\!\left[
K_k^{\mathrm{eff}}\nu_k,
-\delta\theta_{\max},\delta\theta_{\max}
\right],\nonumber\\
\delta\theta_k
&=\delta\theta_{k-1}
+\operatorname{clip}\!\left[
\delta\theta_k^{\star}-\delta\theta_{k-1},
-\dot{\delta\theta}_{\max}T_s,
\dot{\delta\theta}_{\max}T_s
\right],\nonumber\\
\hat\theta_k^{\mathrm{fus}}
&=\hat\theta_k^{\mathrm{MTCN}}+\delta\theta_k,
\label{eq:adaptive_slope_fusion}
\end{align}
where $\delta\theta_{\max}=0.5^{\circ}$ and $\dot{\delta\theta}_{\max}=5^{\circ}\,\mathrm{s}^{-1}$. The same rate-limited update is retained when the target correction changes sign, so the correction crosses zero continuously. If ModernTCN is not ready, the qualified observer is invalid, $|\nu_k|>2^{\circ}$, or $\mathrm{NIS}_k>25$, the correction is reset immediately and $\hat\theta_k^{\mathrm{fus}}=\hat\theta_k^{\mathrm{MTCN}}$ exactly. Re-entry starts from zero and again follows the rate limit. This reset is the only permitted exception to the normal correction-rate constraint.

Before entering the controller, the fused estimate is passed through an amplitude guard,
\begin{equation}
\theta_k^{\mathrm{guard}}=
\begin{cases}
\operatorname{clip}\!\left(
\hat\theta_k^{\mathrm{fus}},-10^{\circ},10^{\circ}
\right),
&|\hat\theta_k^{\mathrm{fus}}|\geq1.3^{\circ},\\
0,
&|\hat\theta_k^{\mathrm{fus}}|<1.3^{\circ}.
\end{cases}
\label{eq:fused_slope_amplitude_guard}
\end{equation}
Only the amplitude deadband and physical clipping are used in this guard; neither the predicted operating-condition label nor its confidence is used to hard-zero an otherwise admissible grade estimate. The guarded signal is subsequently processed by the unchanged causal scheduling filter,
\begin{equation}
\theta_k^{\mathrm{sch}}
=\mathcal F_{\rho}\!\left(\theta_k^{\mathrm{guard}}\right),
\label{eq:fused_slope_conditioning}
\end{equation}
and the resulting $\theta_k^{\mathrm{sch}}$ is the single grade signal supplied to all grade-dependent controller updates. In this way, uncertainty adaptation determines when information from the qualified observer may refine the estimate, while the correction bounds, exact fallback, amplitude guard, and scheduling filter constrain how that refinement can affect the closed loop.

Algorithm~\ref{alg:online_slope_fusion} summarizes the stateful execution order of the fusion interface. The algorithm does not replace the definitions above; it makes explicit the validity branch, exact fallback, correction-state update, and control-side conditioning that must occur in this order at every sampling instant.

\begin{algorithm}[!t]
\caption{Online uncertainty-adaptive grade fusion and control-side scheduling}
\label{alg:online_slope_fusion}
\begin{algorithmic}[1]
\Require conditioned ModernTCN estimate $\hat\theta_k^{\mathrm{MTCN}}$ and confidence $c_k^{\mathrm{main}}$; current six-axis IMU packet $[\tilde{\bm a}_k^{T},\tilde{\bm\omega}_k^{T}]^{T}$; previous RKF, correction, and scheduling-filter states
\Ensure scheduled grade $\theta_k^{\mathrm{sch}}$ and updated persistent states
\State Apply bias compensation and execute the observer prediction using \eqref{eq:qualified_observer_prediction}
\State Adapt the external-acceleration covariance and update the observer using \eqref{eq:qualified_observer_adaptation}
\State Compute $\hat\theta_k^{\mathrm{obs}}$, apply the observer validity check, and determine its disturbance bin using \eqref{eq:rkf_imu_grade_output} and \eqref{eq:fusion_rkf_disturbance_score}
\State Reconstruct $P_{\mathrm{M},k}$, $P_{\mathrm{O},k}$, $C_k$, $S_k$, $\mathrm{NIS}_k$, and $K_k^{\mathrm{eff}}$ using \eqref{eq:moderntcn_regime_uncertainty}--\eqref{eq:innovation_conditioned_imu_gain}
\If{ModernTCN is not ready, the qualified observer is invalid, $|\nu_k|>2^{\circ}$, or $\mathrm{NIS}_k>25$}
    \State $\delta\theta_k\gets0$ and $\hat\theta_k^{\mathrm{fus}}\gets\hat\theta_k^{\mathrm{MTCN}}$ \Comment{exact fallback}
\Else
    \State Compute $\delta\theta_k^{\star}$ and update $\delta\theta_k$ using \eqref{eq:adaptive_slope_fusion}, including rate-limited traversal through correction sign changes
    \State $\hat\theta_k^{\mathrm{fus}}\gets\hat\theta_k^{\mathrm{MTCN}}+\delta\theta_k$
\EndIf
\State Apply \eqref{eq:fused_slope_amplitude_guard} to obtain $\theta_k^{\mathrm{guard}}$
\State Update the causal scheduling filter: $\theta_k^{\mathrm{sch}}\gets\mathcal F_{\rho}(\theta_k^{\mathrm{guard}})$
\State \Return $\theta_k^{\mathrm{sch}}$ and the updated persistent states
\end{algorithmic}
\end{algorithm}

\subsection{Fused-Grade-Scheduled Discrete LPV-MPC}
\label{sec:model_lpv_mpc}

The fusion module becomes control-relevant only when its output modifies the prediction model and optimization problem consistently. The nonlinear plant is therefore converted to a discrete LPV prediction model using the same actuator update and numerical transition as the closed-loop simulation, after which the conditioned grade in \eqref{eq:fused_slope_conditioning} synchronously schedules the model, measured-disturbance path, feedforward term, weights, and bounds, following the established LPV-MPC practice of offline model construction with online scheduled updates \citep{Uchihori2021AUVLPVMPC}.

\subsubsection{Discrete Path-Error Model}

The controller state, input, and measured grade disturbance are
\begin{equation}
\begin{aligned}
\bm{x}^{e}_k
&=\begin{bmatrix}e_{y,k}&e_{\psi,k}&e_{v,k}&e_{\omega,k}^{\mathrm{trk}}\end{bmatrix}^{T},\\
\bm{u}_k
&=\begin{bmatrix}F_{cmd,k}&\omega_{cmd,k}\end{bmatrix}^{T},
\qquad d_k=\theta^{\mathrm{sch}}_k .
\end{aligned}
\label{eq:discrete_error_state}
\end{equation}
Let $\Phi_{T_s}$ denote the implemented one-step transition of the full state $\bm\zeta$ in \eqref{eq:full_nonlinear_state}, including steering-rate and steering-angle saturation, midpoint steering values, fourth-order Runge--Kutta integration, angle wrapping, non-negative-speed protection, and sideslip limiting. The lifting map $\mathcal L_e(\bm x^e;\bm r^{ref},\bar{\bm\zeta}_{int})$ reconstructs the absolute motion state from the path-error coordinates and the local reference pose, speed, and yaw rate, and inserts the steering, sideslip, current, and filter states $\bar{\bm\zeta}_{int}$ associated with the local operating point. Conversely, $\bm q(\bm\zeta;\bm r^{ref})$ maps the propagated absolute state back to the four path errors. The implemented reduced transition is therefore
\begin{equation}
\bm g(\bm x^e,\bm u,\theta)
=\bm q\!\left(
\Phi_{T_s}\!\left[
\mathcal L_e(\bm x^e;\bm r^{ref},\bar{\bm\zeta}_{int}),
\bm u,\theta\right],\bm r^{ref,+}
\right).
\label{eq:implemented_error_transition}
\end{equation}
During nonlinear simulation, all components of $\bm\zeta$ are advanced. During offline database construction, the reduced four-state model fixes the unrepresented internal states at their local operating values before finite differentiation. This makes the reduction and its approximation explicit. Define the operating-point deviations
\begin{equation}
\widetilde{\bm x}^{e}_k=\bm x^e_k-\bar{\bm x}^{e},\qquad
\widetilde{\bm u}_k=\bm u_k-\bar{\bm u},\qquad
\widetilde\theta_k=\theta_k-\bar\theta,
\label{eq:operating_point_deviations}
\end{equation}
whereas the control increment is reserved for the adjacent-sample difference $\Delta\bm u_k=\bm u_k-\bm u_{k-1}$. The local deviation model is
\begin{align}
\widetilde{\bm{x}}^{e}_{k+1}
&=A_d(\boldsymbol\varrho_k)\widetilde{\bm{x}}^{e}_k
+B_d(\boldsymbol\varrho_k)\widetilde{\bm{u}}_k
+E_d(\boldsymbol\varrho_k)\widetilde\theta_k,\nonumber\\
\widetilde{\bm y}_k&=C_d\widetilde{\bm x}^e_k+D_d\widetilde{\bm u}_k,
\qquad
\boldsymbol\varrho_k=
\begin{bmatrix}v_k&\omega_k&\theta^{\mathrm{sch}}_k\end{bmatrix}^{T},
\label{eq:lpv_state}
\end{align}
where $A_d\in\mathbb R^{4\times4}$, $B_d\in\mathbb R^{4\times2}$, $E_d\in\mathbb R^{4\times1}$, $C_d=I_4$, and $D_d=0$. The matrices are evaluated directly from the one-step transition,
\begin{equation}
A_d=\left.\frac{\partial\bm g}{\partial\bm x^e}\right|_{\bar{\boldsymbol\varrho}},
\qquad
B_d=\left.\frac{\partial\bm g}{\partial\bm u}\right|_{\bar{\boldsymbol\varrho}},
\qquad
E_d=\left.\frac{\partial\bm g}{\partial\theta}\right|_{\bar{\boldsymbol\varrho}},
\label{eq:discrete_transition_jacobians}
\end{equation}
using the implemented one-sided state perturbations and centered input and grade perturbations. Thus, no additional continuous-to-discrete conversion is applied.

\subsubsection{Offline LPV Database}

The scheduling grid is formed over speed, yaw rate, and grade. At each grid point, the nominal longitudinal force is
\begin{equation}
\begin{aligned}
F_{cmd}^{\star}
&=c_rmg\cos\theta^{\star}
+\frac{1}{2}\rho_{air}C_DA(v^{\star})^2
+mg\sin\theta^{\star},\\
\bm u^{\star}
&=\begin{bmatrix}F_{cmd}^{\star}&\omega^{\star}\end{bmatrix}^{T}.
\end{aligned}
\label{eq:workpoint_force}
\end{equation}
The database uses 11 speed points from 0.02 to 1.20 m/s, 15 yaw-rate points from $-1.2$ to 1.2 rad/s, and 25 grade points from $-12^{\circ}$ to $12^{\circ}$. It therefore stores $11\times15\times25=4125$ local models, each containing $(A_d,B_d,C_d,D_d,E_d)$, its operating point, and numerical metadata. Online model updating is consequently reduced to interpolation rather than repeated nonlinear linearization.

\subsubsection{Synchronized Scheduling and Interpolation}

The same conditioned scalar is supplied to every grade-dependent controller path,
\begin{equation}
\begin{aligned}
\boldsymbol\varrho_k
&=\begin{bmatrix}v_k&\omega_k&\theta^{\mathrm{sch}}_k\end{bmatrix}^{T},\\
d_k&=\theta^{\mathrm{sch}}_k,\\
F_{eq,k}&=mg\!\left(
\sin\theta^{\mathrm{sch}}_k+c_r\cos\theta^{\mathrm{sch}}_k
\right).
\end{aligned}
\label{eq:synchronized_slope}
\end{equation}
This single-source rule prevents the LPV matrices, measured-disturbance channel, and grade feedforward from receiving differently filtered versions of the grade. The tracking weights, input weights, increment weights, and input envelopes are updated from the same operating condition.

\begin{table*}[!t]
\centering
\caption{\textbf{Implemented LPV-MPC settings and admissible online ranges used in closed-loop simulation.}}
\label{tab:mpc_settings}
\setlength{\tabcolsep}{3pt}
\begin{tabular}{|p{0.19\textwidth}|p{0.13\textwidth}|p{0.34\textwidth}|p{0.25\textwidth}|}
\hline
Item & Symbol & Value & Unit / Description \\
\hline
Sampling time & $T_s$ & 0.0100 & s \\
Prediction horizon & $N_p$ & 150 & 1.50 s \\
Control horizon & $N_c$ & 30 & 0.30 s \\
Nominal tracking weight & $Q_0$ & $\operatorname{diag}(100,100,15,3)$ & $e_y,e_\psi,e_v,e_\omega^{\mathrm{trk}}$ \\
Nominal input weight & $R_0$ & $\operatorname{diag}(3\!\times\!10^{-5},3\!\times\!10^{-5})$ & $F_{cmd},\omega_{cmd}$ \\
Nominal increment weight & $R_{\Delta,0}$ & $\operatorname{diag}(10^{-3},10^{-3})$ & $\Delta F_{cmd},\Delta\omega_{cmd}$ \\
Online range & $Q_k$ & $\operatorname{diag}(50,50,7.5,1.5)$ to $\operatorname{diag}(150,150,22.5,4.5)$ & before condition-dependent gains \\
Online range & $R_k$ & $\operatorname{diag}(1.5,1.5)\!\times\!10^{-5}$ to $\operatorname{diag}(4.5,4.5)\!\times\!10^{-5}$ & scheduled input penalty \\
Online range & $R_{\Delta,k}$ & $\operatorname{diag}(0.5,0.5)\!\times\!10^{-3}$ to $\operatorname{diag}(1.5,1.5)\!\times\!10^{-3}$ & scheduled increment penalty \\
Input constraint & $F_{cmd}$ & $[-600,600]$ & N; online envelope $[-720,720]$ \\
Input constraint & $\omega_{cmd}$ & $[-1.2,1.2]$ & rad/s; online envelope $[-1.44,1.44]$ \\
Increment constraint & $\Delta F_{cmd}$ & $[-400,400]$ & N/step \\
Increment constraint & $\Delta\omega_{cmd}$ & $[-0.9,0.9]$ & (rad/s)/step \\
Soft output constraint & $e_y$ & $[-1.0,1.0]$ & m \\
Soft output constraint & $e_\psi$ & $[-0.5,0.5]$ & rad \\
Runtime slack penalty & $w_{\mathrm{ECR}}$ & $10^4$ & external Adaptive MPC weight \\
Implementation & Solver & MATLAB MPC Toolbox & Adaptive MPC with online model update \\
\hline
\end{tabular}
\end{table*}


After $\boldsymbol\varrho_k$ is saturated to the database limits, trilinear interpolation is performed within the enclosing cell. For any stored matrix $M\in\{A_d,B_d,C_d,D_d,E_d\}$,
\begin{equation}
\begin{aligned}
\xi_r&=\frac{\varrho_{k,r}-\varrho^-_r}{\varrho^+_r-\varrho^-_r},\\
\lambda_{\bm b}&=\prod_{r=1}^{3}\xi_r^{b_r}(1-\xi_r)^{1-b_r},\\
M(\boldsymbol\varrho_k)&=\sum_{\bm b\in\{0,1\}^3}
\lambda_{\bm b}M_{\bm b},
\qquad
\sum_{\bm b}\lambda_{\bm b}=1 .
\end{aligned}
\label{eq:trilinear_interp}
\end{equation}
The interpolated model is held fixed within the current prediction horizon. Smooth condition-dependent maps use $h(z)=3z^2-2z^3$ before scaling configured weights and bounds, avoiding abrupt changes at scheduling-cell boundaries.

\subsubsection{Constrained MPC Optimization}

No road-grade preview is assumed. The current conditioned estimate is held over the prediction horizon, $\theta_{k+i|k}=\theta_k^{\mathrm{sch}}$ for $i=0,\ldots,N_p$, and the corresponding equilibrium command is updated at every sampling instant. The grade disturbance is included in the lifted deviation prediction through
\begin{equation}
\bm\upsilon_k=
\begin{bmatrix}\widetilde{\bm u}_k\\\widetilde\theta_k\end{bmatrix},
\qquad
\bar B_k=
\begin{bmatrix}B_d(\boldsymbol\varrho_k)&E_d(\boldsymbol\varrho_k)\end{bmatrix},
\label{eq:augmented_prediction_model}
\end{equation}
which yields
\begin{equation}
\bm Y_k=\bm F_k\widetilde{\bm x}^e_k
+\bm G_k\widetilde{\bm U}_k
+\bm H_k\widetilde{\bm\Theta}_k .
\label{eq:lifted_prediction}
\end{equation}
The scheduled feedforward term is used to form the input operating point,
\begin{equation}
\bar{\bm u}_k=
\begin{bmatrix}
F_{eq,k}+\tfrac12\rho_{air}C_DA v_k^2 & \omega_k^{ref}
\end{bmatrix}^{\mathsf T}.
\label{eq:online_input_operating_point}
\end{equation}
Thus, $F_{eq,k}$ is not added again after optimization: the optimizer returns the absolute command $\bm u_{k|k}=\bar{\bm u}_k+\widetilde{\bm u}_{k|k}$, and only that command is applied to the plant. The finite-horizon objective is
\begin{equation}
\begin{aligned}
J_k={}&\sum_{i=1}^{N_p}
\|\bm y_{k+i|k}-\bm y^{ref}_{k+i}\|_{Q_k}^{2}\\
&+\sum_{i=0}^{N_c-1}\|\widetilde{\bm u}_{k+i|k}\|_{R_k}^{2}
+\sum_{i=0}^{N_c-1}
\|\Delta\bm u_{k+i|k}\|_{R_{\Delta,k}}^{2}\\
&+w_{\mathrm{ECR}}\|\bm\epsilon_k\|_2^2,
\end{aligned}
\label{eq:mpc_cost}
\end{equation}
subject to
\begin{equation}
\bm u_{\min,k}\leq\bm u_{k+i|k}\leq\bm u_{\max,k},
\qquad
\Delta\bm u_{\min}\leq\Delta\bm u_{k+i|k}
\leq\Delta\bm u_{\max}.
\label{eq:input_constraint}
\end{equation}
where $\Delta\bm u_{k+i|k}=\bm u_{k+i|k}-\bm u_{k+i-1|k}$, with $\bm u_{k-1|k}$ equal to the previously applied command. The lateral and heading errors use the soft bounds
\begin{equation}
\bm y_{\min}-\bm\epsilon_k\leq\bm y_{k+i|k}
\leq\bm y_{\max}+\bm\epsilon_k,
\qquad \bm\epsilon_k\geq\bm0,
\label{eq:soft_output_constraint}
\end{equation}
and the runtime slack penalty $w_{\mathrm{ECR}}$ is held fixed across controller variants.

The controller structure, nominal values, interpolation maps, admissible ranges, and condition-dependent update laws in Table~\ref{tab:mpc_settings} were fixed before comparison. The realized $Q_k$, $R_k$, $R_{\Delta,k}$, and input envelopes can nevertheless vary online because the grade-source signal is the deliberately varied scheduler input and the same fixed update law is applied in every case. Accordingly, a fixed controller configuration means the absence of case-specific retuning, not equality of instantaneous scheduled parameters.

\subsubsection{Online Receding-Horizon Execution}

At every sampling instant, the deployed loop executes the following sequence:
\begin{enumerate}
\item acquire the current AGV responses and form the 22-D feature vector;
\item normalize and augment the features, update the historical window, and evaluate the ModernTCN estimator;
\item execute Algorithm~\ref{alg:online_slope_fusion} to update the qualified observer, fuse the two estimates, and obtain the conditioned grade $\theta_k^{\mathrm{sch}}$;
\item use the same $\theta^{\mathrm{sch}}_k$ to form the scheduling vector, measured-disturbance input, grade feedforward, weights, and bounds;
\item interpolate $(A_d,B_d,C_d,D_d,E_d)$ and solve the constrained quadratic program with $N_p=150$ and $N_c=30$; and
\item apply only the first optimized command to the nonlinear plant and repeat the sequence at the next sample.
\end{enumerate}
This ordered implementation preserves a traceable interface between temporal perception and model-based control while keeping the plant, prediction model, and disturbance update synchronized.

\section{Experimental Setup and Evaluation Protocol}
\label{sec:experiment_design}

\subsection{Research Questions and Evaluation Design}

The evaluation was organized around five sequential research questions. RQ1 asked whether road-grade scheduling materially affects the closed loop of the LPV-MPC with fixed tuning. RQ2 asked whether ModernTCN-delta improves offline grade estimation relative to ModernTCN-22D, GRU-22D, and TCN-22D under a common data protocol. RQ3 asked whether the offline advantage transfers to closed-loop tracking and input regularity. RQ4 asked whether the qualified observer provides complementary information and whether an empirically calibrated, innovation-consistent correction can exploit that information without replacing the learned estimate. RQ5 asked whether the resulting Fusion-LPV-MPC improves matched closed-loop tracking and input regularity relative to MTCN-LPV-MPC, and what core inference cost is measured for the four estimators. This ordering separates the value of the scheduling variable from estimator selection, fusion behavior, and estimator-level computation.

All reported results are computed from data, models, and settings fixed before statistical aggregation. The evaluation distinguishes descriptive means, paired performance effects, acceptance criteria, and deployment claims; a favorable confidence interval for a performance metric does not by itself establish that a separate safety or consistency check is passed.

\subsection{Simulation Platform and Fixed Controller Configuration}

All closed-loop cases used the same revision of the nonlinear AGV plant model and the same LPV-MPC implementation. The controller sample period was 0.01~s, the prediction horizon was 150 steps (1.5~s), and the control horizon was 30 steps (0.3~s). Controller structure, nominal weights, scheduling maps and ranges, increment and output bounds, soft-constraint weights, LPV database, and controller cache were fixed before any comparison was run. The realized scheduled weights and input envelopes were allowed to vary through the same fixed update law because grade source was the experimental factor; no case-specific retuning was performed. Table~\ref{tab:frozen_experimental_configuration} reports the executed configuration rather than historical defaults.

\begin{table*}[!t]
\centering
\caption{\textbf{Experimental and controller configuration held fixed across all reported comparisons.}}
\label{tab:frozen_experimental_configuration}
\scriptsize
\setlength{\tabcolsep}{3.4pt}
\setlength{\tabcolsep}{3pt}
\begin{tabular}{|p{0.20\textwidth}|p{0.23\textwidth}|p{0.49\textwidth}|}
\hline
Item & Fixed value & Execution meaning \\
\hline
Plant and sampling & fixed plant revision (v2); $T_s=0.01$~s & Identical nonlinear plant and sampling interface for every controller case \\
Dataset & 588102 raw samples; 102 continuous runs & 51 route templates with two realizations per template; generator seed 20260511 \\
Reference speed & Training templates: \mbox{$0.80$--$1.18$~m/s} (sample-weighted mean, $0.992$~m/s); closed-loop routes: \mbox{$0.82$--$0.93$~m/s} & Separates the data-generation coverage from the lower-speed closed-loop benchmark \\
History and input & 128 steps (1.28~s); 22D or 88D & The 88D input is $[x_t,\Delta_1x_t,\Delta_2x_t,\Delta_4x_t]$ after train-set normalization; deltas are clipped at 5 \\
Run split & 71/15/16 train/validation/test runs & Run-disjoint; scalers fitted on training runs only; no formal six-route file or sample was reused \\
MPC horizons & $N_p=150$, $N_c=30$ & 1.5-s prediction and 0.3-s control horizons \\
Nominal weights & $Q_0=\mathrm{diag}(100,100,15,3)$; $R_0=\mathrm{diag}(3,3)\!\times\!10^{-5}$; $R_{\Delta,0}=10^{-3}I$ & Same objective form and fixed scheduling law; online factors span 0.5--1.5 before condition-dependent gains \\
Input bounds & Base $[-600,600]$~N and $[-1.2,1.2]$~rad/s & Fixed update law; maximum online envelopes $[-720,720]$~N and $[-1.44,1.44]$~rad/s \\
Increment bounds & $[-400,-0.9]$ to $[400,0.9]$ & Per-step command-increment limits \\
Output bounds & $[-1,-0.5,-0.5,-0.3]$ to $[1,0.5,0.5,0.3]$ & Fixed soft output constraints; position and yaw soft weights were both 3000 \\
Primary estimator & ModernTCN-delta, seeds 1, 7, 11, 21, 42, 73, 101, 202, 340, 520 & Seed 42 was pre-specified for qualitative traces; all ten seeds entered quantitative summaries \\
\hline
\end{tabular}
\end{table*}

\subsection{Simulation Data Generation and Dataset Construction}

The dataset generator used master seed 20260511 to produce 102 continuous simulation runs from 51 route templates, with two random realizations for each template. Windowing was applied only after the run-level split, yielding 16529 training, 3695 validation, and 3602 test windows. The target corresponded to the current window end. Only causal filtering reproducible in online operation was allowed, and normalization statistics were fitted to the training runs before application to validation, test, and online inputs. None of the six formal closed-loop route files or their recorded samples was used to form these windows.

The base input contains 22 proprioceptive and dynamics-derived channels. ModernTCN-delta augments each normalized vector with first differences at lags 1, 2, and 4, producing 88 channels while retaining the same 128-step history. It was trained with AdamW for at most 120 epochs using batch size 256, learning rate $10^{-3}$, weight decay $10^{-4}$, early-stopping patience 25, and a minimum of 30 epochs; the retained checkpoint followed a pre-specified validation rule for each of the ten seeds. Algorithm~\ref{alg:dataset_construction} summarizes the leakage-controlled construction order and makes explicit that run assignment precedes window extraction and that true grade is retained only as an endpoint label.

\begin{algorithm}[!t]
\caption{Leakage-controlled construction of run-disjoint 22D and 88D datasets}
\label{alg:dataset_construction}
\begin{algorithmic}[1]
\Require 51 route templates $\mathcal T$; two fixed realizations per template; causal 22D extractor; $L=128$; $\mathcal L=\{1,2,4\}$
\Ensure run-disjoint paired 22D and ModernTCN-delta 88D window sets $\mathcal D_{\mathrm{tr}}$, $\mathcal D_{\mathrm{va}}$, and $\mathcal D_{\mathrm{te}}$
\State $\mathcal R\gets\varnothing$
\For{each template $\tau\in\mathcal T$ and realization $r\in\{1,2\}$}
    \State Record one continuous run with causal 22D responses, endpoint grade targets, and run/template identifiers; append it to $\mathcal R$
\EndFor
\State Verify $|\mathcal R|=102$ and 588102 total raw samples
\State Split run identifiers 71/15/16 before windowing \Comment{run-disjoint split}
\State Fit $\boldsymbol\mu_{\mathrm{tr}}$ and $\boldsymbol\sigma_{\mathrm{tr}}$ using training-run 22D samples only
\For{each split $s$ and each continuous run $R\in s$}
    \State Normalize $R$ using $\boldsymbol\mu_{\mathrm{tr}}$ and $\boldsymbol\sigma_{\mathrm{tr}}$
    \State Form edge-padded, $\pm5$-clipped $\Delta_{\ell}\mathbf z_k$ for $\ell\in\mathcal L$ and concatenate $\widetilde{\mathbf z}_k=[\mathbf z_k,\Delta_1\mathbf z_k,\Delta_2\mathbf z_k,\Delta_4\mathbf z_k]\in\mathbb R^{88}$
    \State Extract the endpoint-aligned 128-step 22D/88D windows wholly within $R$; attach endpoint grade labels and append to $\mathcal D_s$
\EndFor
\State Verify $|\mathcal D_{\mathrm{tr}}|/|\mathcal D_{\mathrm{va}}|/|\mathcal D_{\mathrm{te}}|=16529/3695/3602$
\State \Return $\mathcal D_{\mathrm{tr}}$, $\mathcal D_{\mathrm{va}}$, and $\mathcal D_{\mathrm{te}}$
\end{algorithmic}
\end{algorithm}

\subsection{Comparison Methods and Grade-Source Controls}

The comparison protocol separated estimator architecture, input representation, scheduling necessity, and fusion value. ModernTCN-22D, GRU-22D, and TCN-22D shared the 22D input and common data protocol, whereas ModernTCN-delta was paired with ModernTCN-22D to isolate the multi-scale delta bank in the offline estimator study. Zero-grade, qualified observer, and truth-driven conditioned sources established the role of grade scheduling, while the final closed-loop benchmark used the repaired MTCN-LPV-MPC interface and its matched Fusion-LPV-MPC counterpart. The latter pair isolated the bounded correction from the qualified observer; the truth-driven source remained an analysis-only reference. Across controller comparisons, the plant, initialization, LPV database, controller structure, nominal settings, scheduling maps, admissible ranges, and update laws were fixed without case-specific retuning. Realized scheduled weights and input envelopes could nevertheless vary through the same fixed update law because grade source was the experimental factor. Five-source means are reported only as descriptive context because deterministic and learned sources retain different statistical units.

\subsection{Closed-Loop Routes, Seeds, and Statistical Units}

The formal controller benchmark used six predefined routes: P1 Factory logistics, P2 Sharp-turn transition, P3 Long up/down slope, P4 Mild slope-turn coupling, P5 Flat factory logistics, and P6 Downhill recovery. The deterministic zero-grade, qualified observer, and truth-driven sources contributed one case per route. Learned sources contributed ten model-seed cases per route. Consequently, deterministic comparisons were paired by route, whereas learned comparisons preserved the model-seed structure before route aggregation. Figure~\ref{fig:closed_loop_route_set} shows the six evaluation routes.

The fusion calibration and configuration selection used only the recorded training and validation runs. The empirical uncertainty of the qualified observer was estimated from the 71 training runs, while ModernTCN uncertainty, cross-source error dependence, candidate validation, and candidate ranking used the 15 validation runs of the ten trained models. The final closed-loop Fusion comparison used the same six routes and ten model seeds as MTCN-LPV-MPC, giving 60 matched route--seed pairs. Quantitative conclusions use all 60 pairs; seed-level records are retained as supplementary data rather than enumerated in the main text.

\begin{figure*}[!t]
\centering
\includegraphics[width=\textwidth]{../7.6/figures_final/Fig03_closed_loop_route_set/output/fig03_closed_loop_route_set.pdf}
\caption{Six-route closed-loop benchmark. Each panel shows the complete planar reference path together with the aligned grade and reference-speed profiles. Green circles, red squares, and blue arrowheads mark route starts, route ends, and travel direction, respectively.}
\label{fig:closed_loop_route_set}
\end{figure*}

Model seed 42 was fixed in advance for the qualitative estimator and fusion traces. Quantitative conclusions use all ten seeds and all six formal routes; qualitative traces illustrate the runtime mechanism and do not replace the paired summaries.

\subsection{Metrics and Statistical Analysis}

Offline primary metrics were grade mean absolute error (MAE) and either the pre-specified $|\theta|\leq10^{\circ}$ P95 absolute error or the edge-region P95 absolute error, depending on the comparison. The frozen edge-region metric was the larger of the P95 absolute errors evaluated over $-10^{\circ}\leq\theta\leq-8^{\circ}$ and $8^{\circ}\leq\theta\leq10^{\circ}$. The three inferential primary closed-loop outcomes were lateral-error RMS, heading-error RMS, and the composite input-increment index $J_{\Delta u}$. Position-error RMS and yaw-rate-command RMS were retained as secondary descriptive outcomes and are reported for completeness. Lower values favor the target method. The composite input-increment index is a control-regularity summary, not a physical energy measure.

Paired effects are reported as comparator minus target so that positive values favor the target. Offline intervals resampled paired model seeds. Learned closed-loop intervals resampled model seeds first and then retained all six matched routes within each resampled seed; deterministic comparisons resampled routes. All reported bootstrap intervals used 10000 iterations and seed 20260715. Table~\ref{tab:evaluation_protocol} records the unit and interpretation of each evidence family.

\begin{table*}[!t]
\centering
\caption{\textbf{Evaluation metrics, statistical units, and interpretation rules.}}
\label{tab:evaluation_protocol}
\scriptsize
\setlength{\tabcolsep}{3.0pt}
\setlength{\tabcolsep}{3pt}
\begin{tabular}{|p{0.19\textwidth}|p{0.25\textwidth}|p{0.16\textwidth}|p{0.16\textwidth}|p{0.16\textwidth}|}
\hline
Evidence family & Reported quantity & Unit & Statistical unit & Decision rule \\
\hline
Offline grade estimation & MAE; pre-specified P95 absolute error; edge-region P95 absolute error & deg & Ten paired model seeds & Positive comparator-minus-target effect favors target \\
Primary closed-loop inference & $e_y$ RMS; $e_\psi$ RMS; $J_{\Delta u}$ & m; rad; -- & Six paired routes or ten seeds $\times$ six routes & Route or hierarchical seed-first bootstrap \\
Secondary closed-loop description & Position-error RMS; $\omega_{cmd}$ RMS & m; rad/s & Same cases as the associated comparison & Descriptive complement to the primary metrics \\
Fusion calibration & Validation MAE ratio; active-sample fractions; $K^{\mathrm{raw}}$/$K^{\mathrm{eff}}$ spread and cap fraction & ratio; fraction & 15 validation runs $\times$ ten model seeds & Pre-specified acceptance criteria followed by the candidate ranking \\
Fusion closed-loop comparison & Scheduled-grade MAE; $e_y$ RMS; $e_\psi$ RMS; $J_{\Delta u}$ & deg; m; rad; -- & Ten seeds $\times$ six matched routes & Descriptive source accuracy plus hierarchical seed-first paired bootstrap for the three primary closed-loop metrics \\
Safety checks & Pass, fail, or not assessable & status & Pre-specified comparison/case & Missing inputs or zero denominators are never coerced to pass \\
Estimator core runtime & p50, p95, p99, $>10$-ms count & ms & 10000 timed calls after 500 warm-ups & Compare methods only within the common ONNX Runtime CPU backend \\
\hline
\end{tabular}
\end{table*}

\subsection{Reproducibility and Reporting Boundaries}

The dataset, feature specification, model checkpoints, plant revision, six route definitions, LPV database, controller cache, metrics, thresholds, and seed set were fixed and identified by file checksums. A completeness check confirmed 40/40 offline estimator rows, 240/240 estimator closed-loop rows, 18/18 scheduling-necessity cases, and 138/138 five-source controller cases. A source cross-check covered 38 confidence intervals with no discrepancy.

The protocol includes no physical-vehicle or hardware-in-the-loop test and no formal stability proof.

\section{Results and Discussion}
\label{sec:results_discussion}

\subsection{RQ1: Does Accurate Grade Scheduling Affect the LPV-MPC Closed Loop?}

\begin{figure*}[!t]
\centering
\includegraphics[width=\textwidth]{../7.6/figures_final_1/Fig04_slope_necessity_effects/output/fig04_slope_necessity_effects_ra_x2_qg_x05.pdf}
\caption{Route-resolved absolute closed-loop metrics under zero-grade scheduling, the truth-driven conditioned reference, and qualified-observer scheduling. Panels show (a) lateral-error RMS, (b) heading-error RMS, and (c) the composite input-increment index. White circles denote zero-grade scheduling, orange diamonds the truth-driven reference, and teal squares the qualified observer; connectors pair each grade-informed result with the same route's zero-grade comparator. Lower values indicate better performance, and logarithmic axes accommodate route-to-route scale differences. The truth-driven conditioned reference is analysis-only and represents an idealized scheduling reference rather than a deployable grade source.}
\label{fig:slope_necessity_effects}
\end{figure*}


Across the six routes, the absolute means followed the expected aggregate ordering from missing to estimated to exact grade information. Zero-grade scheduling yielded lateral-error RMS, heading-error RMS, and $J_{\Delta u}$ means of 1.48046~m, 0.34204~rad, and 696.014, respectively. The qualified observer reduced these values to 0.02051~m, 0.02347~rad, and 0.51646, while the analysis-only truth-driven reference yielded the lowest corresponding means of 0.01248~m, 0.01139~rad, and 0.42684. Relative to zero-grade scheduling, the mean reductions under the truth-driven reference were 1.46798~m, 0.33065~rad, and 695.587; the corresponding reductions under the qualified observer were 1.45995~m, 0.31857~rad, and 695.497. All six route-bootstrap 95\% intervals excluded zero (Table~\ref{tab:slope_necessity_effects}). These aggregate results support both the control relevance of introducing grade scheduling and the additional benefit of a more accurate scheduling signal.

The route-resolved view in Fig.~\ref{fig:slope_necessity_effects} also shows that this aggregate ordering is not a pointwise monotonic rule. The qualified observer slightly increased lateral-error RMS on P4 and P5 relative to zero-grade scheduling, heading-error RMS on P5, and $J_{\Delta u}$ on P4--P6. Thus, accurate grade scheduling improves the closed loop overall under the tested protocol, whereas the realized benefit of the deployable qualified observer remains dependent on route geometry, vehicle dynamics, and the selected closed-loop metric.

\begin{table*}[!t]
\centering
\caption{\textbf{Six-route paired effects relative to zero-grade LPV-MPC. Effects are zero-grade minus target, so positive values favor the target. Boldface mean effects have 95\% confidence intervals excluding zero. The truth-driven conditioned reference is analysis-only.}}
\label{tab:slope_necessity_effects}
\setlength{\tabcolsep}{3pt}
\begin{tabular}{|l|l|r|r|r|}
\hline
Target source & Metric & Mean effect & 95\% CI & Routes \\
\hline
Truth-driven conditioned reference & $e_y$ RMS (m) & \textbf{1.46798} & [0.11715, 3.74682] & 6 \\
Truth-driven conditioned reference & $e_\psi$ RMS (rad) & \textbf{0.33065} & [0.05065, 0.80004] & 6 \\
Truth-driven conditioned reference & $J_{\Delta u}$ & \textbf{695.587} & [116.408, 1339.454] & 6 \\
Qualified observer & $e_y$ RMS (m) & \textbf{1.45995} & [0.11917, 3.72848] & 6 \\
Qualified observer & $e_\psi$ RMS (rad) & \textbf{0.31857} & [0.04460, 0.79161] & 6 \\
Qualified observer & $J_{\Delta u}$ & \textbf{695.497} & [116.277, 1339.412] & 6 \\
\hline
\end{tabular}
\end{table*}

Figure~\ref{fig:slope_necessity_effects} reports the absolute route-level outcomes and exposes the large differences in operating scale, while Table~\ref{tab:slope_necessity_effects} reports comparator-minus-target mean effects and the route-bootstrap uncertainty used for inference. The route-dependent behavior of the qualified observer is consistent with the residual sensitivity of inertial grade estimation to vehicle acceleration and pitch dynamics, even when external-acceleration compensation is applied \citep{Liu2023RoadSlopeAAIMM}.

\subsection{RQ2: Does the Multi-Scale Delta Bank Improve Offline Grade Estimation?}

\begin{figure*}[!t]
\centering
\includegraphics[width=\textwidth]{../7.6/figures_final/Fig05_offline_estimator_distribution/output/fig05_offline_estimator_distribution.pdf}
\caption{Matched-seed offline comparison of ModernTCN input representations. (a) Grade MAE over test windows with $|\theta|\leq10^{\circ}$. (b) Edge-region P95 absolute grade error. Blue circles and vermillion squares show individual results for ModernTCN-delta and ModernTCN-22D, respectively; gray lines connect matched model seeds, and larger markers denote arithmetic means. Green diamonds show seed-paired effects defined as ModernTCN-22D minus ModernTCN-delta. The larger diamond and vertical whiskers denote the mean effect and its 95\% paired-seed bootstrap confidence interval. Positive effects favor ModernTCN-delta ($n=10$).}
\label{fig:offline_estimator_distribution}
\end{figure*}


Figure~\ref{fig:offline_estimator_distribution} isolates the effect of the multi-scale delta bank using ten matched model seeds. Over test windows with $|\theta|\leq10^{\circ}$, the mean grade MAE decreased from 0.65300$^{\circ}$ for ModernTCN-22D to 0.59888$^{\circ}$ for ModernTCN-delta. The paired effect was 0.05412$^{\circ}$ with a 95\% confidence interval of [0.00608, 0.10180]. In the edge region, the mean P95 absolute error decreased from 2.88256$^{\circ}$ to 2.45649$^{\circ}$, corresponding to a paired effect of 0.42608$^{\circ}$ [0.18517, 0.70987].

The broader four-estimator benchmark is reported in Table~\ref{tab:estimator_and_delta_results}. ModernTCN-delta achieved the lowest mean MAE and pre-specified $|\theta|\leq10^{\circ}$ P95 absolute error, with supported paired advantages over GRU-22D and TCN-22D.

\begin{table*}[!t]
\centering
\caption{\textbf{Offline estimator results. Positive comparator-minus-target effects favor ModernTCN-delta; boldface effects have 95\% confidence intervals excluding zero. Final closed-loop controller means are reported in the matched A4 benchmark.}}
\label{tab:estimator_and_delta_results}
\scriptsize
\setlength{\tabcolsep}{3.0pt}
\setlength{\tabcolsep}{3pt}
\begin{tabular}{|p{0.18\textwidth}|p{0.20\textwidth}|p{0.16\textwidth}|p{0.24\textwidth}|p{0.12\textwidth}|}
\hline
Evidence & Method or comparison & Metric / mean & Effect and 95\% CI & Paired support \\
\hline
Offline mean ($n=10$) & ModernTCN-delta & MAE / P95 ($|\theta|\leq10^{\circ}$): 0.59888 / 1.67777 deg & -- & -- \\
Offline mean ($n=10$) & ModernTCN-22D & MAE / P95 ($|\theta|\leq10^{\circ}$): 0.65300 / 1.80626 deg & -- & -- \\
Offline mean ($n=10$) & GRU-22D & MAE / P95 ($|\theta|\leq10^{\circ}$): 0.88583 / 2.25950 deg & -- & -- \\
Offline mean ($n=10$) & TCN-22D & MAE / P95 ($|\theta|\leq10^{\circ}$): 0.91537 / 2.63580 deg & -- & -- \\
\hline
Representation & Delta vs. ModernTCN-22D & MAE ($|\theta|\leq10^{\circ}$) / edge-region P95 & \textbf{0.05412 [0.00608, 0.10180] / 0.42608 [0.18517, 0.70987] deg} & 8/10; 9/10 seeds \\
Estimator & Delta vs. GRU-22D & MAE / P95 ($|\theta|\leq10^{\circ}$) & \textbf{0.28694 [0.18468, 0.37271] / 0.58172 [0.38462, 0.80045] deg} & 9/10; 10/10 seeds \\
Estimator & Delta vs. TCN-22D & MAE / P95 ($|\theta|\leq10^{\circ}$) & \textbf{0.31648 [0.24668, 0.38012] / 0.95802 [0.69170, 1.25456] deg} & 10/10; 10/10 seeds \\
\hline
\end{tabular}
\end{table*}

Figure~\ref{fig:offline_estimator_distribution} shows all ten matched seed pairs for the two ModernTCN input representations, while Table~\ref{tab:estimator_and_delta_results} retains the complete four-estimator comparison. The evidence supports the tested ModernTCN-delta configuration under the present evaluation protocol; it is not a universal ranking of temporal architectures.

\subsection{RQ3: Does Offline Improvement Transfer Selectively to Closed-Loop Control?}

The final closed-loop MTCN-LPV-MPC benchmark achieved mean lateral-error RMS of 0.02635~m, heading-error RMS of 0.02216~rad, and $J_{\Delta u}=0.93053$ over six routes and ten model seeds. These values are reported as the final controller baseline; the deterministic qualified observer and truth-driven sources remain descriptive route-level references because they do not share the learned-source seed structure. This separation avoids treating estimator-level offline improvements as automatically equivalent to paired closed-loop controller effects: the scheduling interface, route geometry, constrained optimizer, and command regularization transform grade-estimation errors before they appear in tracking metrics.

\subsection{RQ4: Is the Qualified Observer a Reliable Complementary Source?}

\begin{figure}[!tb]
\centering
\includegraphics[width=\columnwidth]{../7.6/figures_final_1/Fig06_imu_observer_trace/output/fig06_imu_observer_trace.pdf}
\caption{Complete P2 Sharp-turn transition traces obtained from the same closed-loop simulation. Panel (a) compares the reference grade with the raw estimates from ModernTCN-delta and the Qualified observer. ModernTCN-delta uses model seed 42, while both estimator branches share the same sensor realization generated with sensor seed 4631. Panel (b) shows the paired samplewise absolute-error difference $|e_{\mathrm{obs}}|-|e_{\mathrm{MTCN}\text{-}\delta}|$; negative values favor the Qualified observer, whereas positive values favor ModernTCN-delta.}
\label{fig:imu_observer_trace}
\end{figure}


The qualified observer evaluated here is a causal RKF-IMU estimator that uses the six-axis accelerometer and gyroscope measurements exclusively. Across the six deterministic closed-loop routes used in RQ1, it achieved a mean scheduled-grade MAE of 1.18210$^{\circ}$. All 62582 logged observer updates were valid, and no constraint violation, solver failure, or timeout was recorded. The same observer is used as the standalone scheduling source in RQ1; fusion-specific calibration and inference are considered separately below.

Figure~\ref{fig:imu_observer_trace} shows the complete 5201-sample, 52-s P2 Sharp-turn transition record obtained by evaluating both estimator branches within the same closed-loop simulation. After excluding the initial 0.5~s, the raw grade-estimation MAE was 1.00461$^{\circ}$ for the qualified observer and 0.62075$^{\circ}$ for ModernTCN-delta. With a 0.1$^{\circ}$ tolerance for a decisive local difference, the qualified observer had lower absolute error on 34.07\% of the evaluated samples, ModernTCN-delta had lower error on 55.95\%, and the remaining 9.98\% were near ties.

The qualified observer evaluated here is a causal RKF-IMU estimator that uses the six-axis accelerometer and gyroscope measurements exclusively. Across the six deterministic closed-loop routes used in RQ1, it achieved a mean scheduled-grade MAE of 1.18210$^{\circ}$. All 62582 logged observer updates were valid, and no constraint violation, solver failure, or timeout was recorded. The same observer is used as the standalone scheduling source in RQ1; fusion-specific calibration and inference are considered separately below.

\begin{figure*}[!t]
\centering
\includegraphics[width=\textwidth]{../7.6/figures_final/Fig07_node44r1_gain_calibration/output/fig07_node44r1_gain_calibration.pdf}
\caption{Validation-only calibration of the uncertainty-adaptive fusion gain. The panels summarize the empirical source-error moments, the distributions of $K^{\mathrm{raw}}$ and $K^{\mathrm{eff}}$, and the candidate validation outcome. $K^{\mathrm{raw}}$ and $K^{\mathrm{eff}}$ denote the unconstrained and the innovation-weighted effective fusion gains.}
\label{fig:fusion_case_distribution}
\end{figure*}

\begin{table*}[!t]
\centering
\caption{\textbf{Qualified observer validation results and fusion gain calibration on the validation runs. Ratios below one favor the fused source.}}
\label{tab:observer_fusion_evidence}
\scriptsize
\setlength{\tabcolsep}{3.2pt}
\setlength{\tabcolsep}{3pt}
\begin{tabular}{|p{0.19\textwidth}|p{0.14\textwidth}|p{0.36\textwidth}|p{0.22\textwidth}|}
\hline
Stage & Cases/pairs & Result & Interpretation \\
\hline
Qualified observer validation & 90 cases & MAE 0.19985$^{\circ}$; P95 0.49109$^{\circ}$; bias $-0.01079^{\circ}$; invalid 1.75\%; safe-beneficial 44.20\% & Qualified observer criteria passed \\
$K_{\max}=0.3$ & 150 validation run--seed cases & MAE ratio 0.92260; $K^{\mathrm{eff}}$ spread 0.01872 & Did not meet the gain-spread criterion \\
$K_{\max}=0.4$ & 150 validation run--seed cases & MAE ratio 0.90302; active 37.23\%; cap 91.95\% & Passed \\
$K_{\max}=0.5$ & 150 validation run--seed cases & MAE ratio 0.88970; active 38.34\%; cap 64.49\% & Passed and selected by the ranking rule \\
Selected-gain verification & 458327 valid samples & $K^{\mathrm{raw}}$ P05/P95 0.28128/0.70826; $K^{\mathrm{eff}}$ P05/P95 0.28128/0.50000 & Condition-dependent gain verified \\
\hline
\end{tabular}
\end{table*}

\begin{figure*}[!t]
\centering
\includegraphics[width=\textwidth]{../7.6/figures_final/Fig08_node44r1_mechanism_trace/output/fig08_node44r1_mechanism_trace.pdf}
\caption{Complete fusion mechanism trace for P2 Sharp-turn transition at model seed 42, which was fixed in advance for qualitative visualization. The panels show the reference, ModernTCN-delta, Qualified observer, and fused grades; the condition-dependent effective gain; the target and applied corrections with rate-limit activity; and the normalized innovation squared (NIS) with fallback intervals.}
\label{fig:fusion_qualitative_trace}
\end{figure*}

\subsubsection{Validation-only calibration of the adaptive gain}

The validation-only calibration evaluated three gain caps: 0.3, 0.4, and 0.5. The 0.3 candidate did not satisfy the pre-specified $K^{\mathrm{eff}}$-spread criterion, whereas 0.4 and 0.5 passed. The selection rule chose $K_{\max}=0.5$, which achieved a macro run--seed MAE ratio of 0.88970, an active fraction of 38.34\%, and active-sample improvement and degradation fractions of 75.88\% and 4.48\%, respectively.

For the selected candidate, the validation P05--P95 intervals were 0.28128--0.70826 for $K^{\mathrm{raw}}$ and 0.28128--0.50000 for $K^{\mathrm{eff}}$. The cap was active for 64.49\% of valid non-fallback samples, while 35.37\% were at least 0.01 below the cap. Thus, 0.5 is an upper bound rather than a fixed fusion weight; the deployed correction varies with the calibrated source errors, confidence, excitation, and innovation consistency (Table~\ref{tab:observer_fusion_evidence}).

Table~\ref{tab:observer_fusion_evidence} summarizes the validation calibration and the non-degenerate distributions of $K^{\mathrm{raw}}$ and $K^{\mathrm{eff}}$. Figure~\ref{fig:fusion_qualitative_trace} then shows the complete P2 Sharp-turn transition trace at model seed 42, relating the source estimates, effective gain, correction state, NIS, and exact-fallback intervals. The trace illustrates the mechanism under a seed fixed before final aggregation and is not used as a substitute for the quantitative summaries.

The validation evidence establishes that $K_{\max}=0.5$ is a cap rather than a constant correction weight. The bounded channel therefore preserves the learned source as the primary estimate while allowing the qualified observer to contribute according to the calibrated error structure and online innovation consistency. This construction is consistent with the broader principle of combining physical and learned information without converting the learning block into the control law \citep{Hewing2020CautiousGPMPC,Zheng2024HybridPhysicsLearningUSV}.

\begin{table*}[!t]
\centering
\caption{\textbf{Descriptive five-source context and matched MTCN-LPV-MPC minus Fusion-LPV-MPC effects. Panel A is not a unified ranking; positive effects favor Fusion in Panel B, and boldface effects have 95\% confidence intervals excluding zero.}}
\label{tab:five_source_closed_loop}
\scriptsize
\setlength{\tabcolsep}{2.7pt}
\setlength{\tabcolsep}{3pt}
\begin{tabular}{|l|r|r|r|r|r|r|}
\hline
\multicolumn{7}{@{}l}{\textit{Panel A: descriptive raw means under source-specific statistical units}} \\
\hline
Grade source & Cases & $e_y$ RMS (m) & $e_\psi$ RMS (rad) & Position RMS (m) & $J_{\Delta u}$ & $\omega_{cmd}$ RMS (rad/s) \\
\hline
Zero-grade LPV-MPC & 6 & 1.48046 & 0.34204 & 9.88861 & 696.014 & 0.38514 \\
Qualified observer & 6 & 0.02051 & 0.02347 & 1.03259 & 0.51646 & 0.05469 \\
MTCN-LPV-MPC & 60 & 0.02635 & 0.02216 & 0.50588 & 0.93053 & 0.05868 \\
Fusion-LPV-MPC & 60 & 0.02301 & 0.02049 & 0.48115 & 0.80093 & 0.05712 \\
Truth-driven conditioned reference & 6 & 0.01248 & 0.01139 & 0.25546 & 0.42684 & 0.05156 \\
\hline
\multicolumn{7}{@{}l}{\textit{Panel B: matched MTCN-minus-Fusion inference}} \\
\hline
Paired metric & Pairs & \multicolumn{2}{l}{Mean effect [95\% CI]} & \multicolumn{2}{l}{Improved/degraded/tied} \\
$e_y$ RMS (m) & 60 & \multicolumn{2}{l}{0.00334 [$-0.00011$, 0.00794]} & \multicolumn{2}{l}{28/14/18} \\
$e_\psi$ RMS (rad) & 60 & \multicolumn{2}{l}{\textbf{0.00167 [0.00026, 0.00334]}} & \multicolumn{2}{l}{33/9/18} \\
$J_{\Delta u}$ & 60 & \multicolumn{2}{l}{\textbf{0.12960 [0.02163, 0.34511]}} & \multicolumn{2}{l}{25/17/18} \\
\hline
\end{tabular}
\end{table*}

\subsection{RQ5: Does Fusion Add Value beyond ModernTCN in Closed Loop?}

\begin{figure*}[!t]
\centering
\includegraphics[width=\textwidth]{../7.6/figures_final/Fig09_node44r1_formal_summary/output/fig09_node44r1_formal_summary.pdf}
\caption{Formal matched Fusion-LPV-MPC versus MTCN-LPV-MPC results over six routes and ten model seeds. The first panel reports hierarchical paired-bootstrap effects and 95\% confidence intervals for lateral-error RMS, heading-error RMS, and the composite input-increment index $J_{\Delta u}$; positive MTCN-minus-Fusion effects favor Fusion. The second panel reports route-level mean effects for the same metrics.}
\label{fig:five_controller_qualitative}
\end{figure*}

The five-source means in Table~\ref{tab:five_source_closed_loop} provide descriptive controller context rather than a unified ranking. Zero-grade, qualified observer, and truth-driven sources contain six route cases each, whereas MTCN-LPV-MPC and Fusion-LPV-MPC contain 60 route--seed cases each. Formal RQ5 inference is limited to the 60 matched MTCN--Fusion route--seed pairs.

The descriptive means indicate complementary strengths rather than a uniform ranking. The qualified observer yielded lower lateral-error RMS, input-increment index, and steering-rate-command RMS, whereas MTCN-LPV-MPC retained lower heading-error and position RMS. Because the qualified observer contributes six deterministic route cases and MTCN-LPV-MPC contributes 60 route--seed cases, these differences remain descriptive; formal RQ5 inference is restricted to the matched MTCN--Fusion pairs.

The fused scheduled-grade MAE decreased from 0.50260$^{\circ}$ to 0.41007$^{\circ}$, a descriptive mean reduction of 0.09253$^{\circ}$ (18.4\%). Route-level grade MAE improved on five routes and was unchanged on the sixth; among the 60 pairs, 40 improved, 18 were unchanged, and two degraded. Because scheduled-grade MAE was retained as descriptive source-accuracy evidence, it is not used to alter the pre-specified inference on the three primary closed-loop metrics.

Fusion reduced the mean of each primary closed-loop metric. The paired MTCN-minus-Fusion effects were 0.00334~m for lateral-error RMS (95\% CI [$-0.00011$, 0.00794]), 0.00167~rad for heading-error RMS (95\% CI [0.00026, 0.00334]), and 0.12960 for $J_{\Delta u}$ (95\% CI [0.02163, 0.34511]). The latter two intervals support improved heading tracking and input regularity; the lateral-error mean was favorable but remained statistically unresolved. No Fusion safety failure was recorded.

The practical relevance of these effects is best judged against the
attainable headroom rather than their absolute size. At the considered
operating points, MTCN-LPV-MPC already tracks with a mean lateral-error
RMS of 0.02635~m, so the remaining improvement potential is bounded by
the truth-driven reference (0.01248~m). Computed from the means in
Table~\ref{tab:five_source_closed_loop}, the bounded correction recovers
24.1\% of this residual lateral-error gap, 15.5\% of the heading-error
gap, and 25.7\% of the input-increment gap, corresponding to relative
reductions of 12.7\%, 7.5\%, and 13.9\% over MTCN-LPV-MPC. A
software-level interface change that closes about one quarter of the gap
toward truth-driven scheduling is therefore a meaningful refinement at
this error scale.

Figure~\ref{fig:five_controller_qualitative} summarizes the aggregate paired effects and their route-level distribution. It combines the three hierarchical bootstrap intervals with all 18 route--metric mean comparisons. Route-level means favored Fusion in 14 of the 18 comparisons and were unchanged in three; the single exception, $J_{\Delta u}$ on P6 Downhill recovery, reflects a localized route sensitivity and does not alter the supported aggregate reduction. Complete route- and seed-level results are retained in the supplementary paired-results data.

\subsection{What Is the Core Inference Cost of the Tested Estimators?}

\begin{table*}[!t]
\centering
\caption{\textbf{Same-backend core inference time and model size for the four tested estimators. p50/p95/p99 denote latency percentiles over 10000 timed calls.}}
\label{tab:runtime_scope}
\scriptsize
\setlength{\tabcolsep}{2.8pt}
\setlength{\tabcolsep}{3pt}
\begin{tabular}{|l|r|r|r|r|r|l|}
\hline
Method & Params (k) & Model (kB) & p50 (ms) & p95 (ms) & p99 (ms) & $>10$ ms \\
\hline
GRU-22D & 99.276 & 400.682 & 0.3071 & 0.3377 & 0.4124 & 0/10000 \\
ModernTCN-22D & 118.538 & 472.135 & 0.3582 & 0.3881 & 0.4942 & 0/10000 \\
ModernTCN-delta & 145.202 & 578.793 & 0.3744 & 0.4038 & 0.4818 & 0/10000 \\
TCN-22D & 156.396 & 627.487 & 1.0452 & 1.1341 & 1.3597 & 1/10000 \\
\hline
\end{tabular}
\end{table*}


Under the common ONNX Runtime CPU backend, batch size one, one thread, 500 warm-ups, and 10000 timed calls, all four estimators had p95 core inference times below 1.2~ms (Table~\ref{tab:runtime_scope}). ModernTCN-delta achieved a p95 of 0.4038~ms with no call above 10~ms. GRU-22D was the fastest core predictor, while TCN-22D produced one call above 10~ms. The sub-millisecond cores leave substantial margin within the 10-ms control period without requiring solver-side acceleration of the kind pursued by parallel NMPC implementations \citep{Xu2024ParallelNMPC}. These same-backend measurements characterize estimator inference only and are not used to rank unmeasured wrappers or controller interfaces.

\subsection{Integrated Discussion and Scope}

Taken together, the experiments reveal a staged rather than monotonic perception-to-control relationship. The truth-driven comparison first establishes that grade is a consequential scheduling variable for the LPV-MPC with fixed tuning. ModernTCN-delta then improves grade estimation under the common offline protocol, but the representation gain transfers most clearly to input regularity rather than uniformly to both tracking errors. This distinction is expected in a constrained scheduled controller: route geometry, state evolution, saturation margins, and command penalties transform local estimation errors before they appear in closed-loop metrics. Offline grade error should therefore support estimator selection, not substitute for controller evaluation.

The qualified observer is useful because its error dynamics differ from those of the learned history-based estimate, but it remains a secondary source. Empirical total-error calibration, cross-error accounting, and innovation conditioning convert this complementarity into a non-degenerate bounded correction: the selected effective gain remained below its cap for a substantial fraction of validation samples, and formal scheduled-grade MAE improved on five routes without changing on the sixth. In closed loop, the aggregate evidence supported lower heading-error RMS and input increments, together with a favorable mean lateral-error reduction. In headroom terms, the bounded correction closes roughly one quarter of the residual gap toward truth-driven scheduling. This pattern aligns with prior use of model-based information to manage grade-estimation uncertainty \citep{Liu2023RoadSlopeAAIMM,Gao2024RoadSlopeHDV} and with hybrid physics--learning designs that retain the model-based controller as the command-generating mechanism \citep{Zheng2024HybridPhysicsLearningUSV}.

\section{Conclusion}
\label{sec:conclusion}

This study developed a slope-aware perception-to-control chain for a diagonal dual-steering-wheel AGV. ModernTCN-delta estimates grade from causal proprioceptive history, the qualified observer supplies bounded complementary information, and Fusion-LPV-MPC uses the conditioned estimate to schedule the model-based controller. Both auxiliary classification heads support representation learning; only the main-condition confidence calibrates fusion uncertainty, while the steering-direction output is retained for diagnostics. Neither predicted label directly switches the controller.

Under the fixed closed-loop protocol, the final MTCN-LPV-MPC benchmark
comprised ten model seeds across six predefined routes, whereas the
qualified observer contributed one deterministic case per route.
Published inertial-sensor-based methods for grade or pitch estimation provide
the relevant literature context
\citep{Jauch2018IMUGradeOrientation,Li2022QuaternionUKFGrade,
Kim2023InertialRoadPitch}; however, none of these published implementations
was reproduced as an experimental baseline. Relative to the qualified observer
implemented in this study, MTCN-LPV-MPC employed the AGV-adapted
ModernTCN-delta estimator. Its temporal backbone was simplified from the
established ModernTCN architecture \citep{Luo2024ModernTCN}. MTCN-LPV-MPC
achieved a mean lateral-error RMS of 0.02635~m, a heading-error RMS of
0.02216~rad, and $J_{\Delta u}=0.93053$. The corresponding qualified observer
route means were 0.02051~m, 0.02347~rad, and 0.51646, respectively. Numerically,
the qualified observer yielded lower lateral-error RMS and input increments,
whereas MTCN-LPV-MPC retained a slightly lower heading-error RMS. Because the
two sources have different statistical units, these results constitute a
descriptive comparison rather than a paired inferential effect. In the
separate ten-seed offline comparison, ModernTCN-delta achieved a grade MAE of
0.59888$^{\circ}$ and a P95 error of 1.67777$^{\circ}$. ModernTCN-22D yielded
0.65300$^{\circ}$ and 1.80626$^{\circ}$, respectively. The corresponding
values were 0.88583$^{\circ}$ and 2.25950$^{\circ}$ for GRU-22D
\citep{Cho2014GRU}, and 0.91537$^{\circ}$ and 2.63580$^{\circ}$ for TCN-22D
\citep{Bai2018TCN,Cao2021TCNRSA}.

Building on the task-specific ModernTCN-delta estimator, whose temporal
backbone was simplified from the established ModernTCN architecture
\citep{Luo2024ModernTCN}, we evaluated the incremental effect of
uncertainty-adaptive fusion. The matched comparison isolated the fusion
interface: MTCN-LPV-MPC used the unfused ModernTCN-delta estimate,
whereas Fusion-LPV-MPC added the calibrated, bounded correction.
Relative to this matched baseline, the fusion interface reduced mean
scheduled-grade MAE from $0.50260^{\circ}$ to $0.41007^{\circ}$,
corresponding to an 18.4\% reduction. Route-level grade MAE decreased
on five routes and remained unchanged on the sixth. Across the 60
matched route--seed cases, Fusion-LPV-MPC lowered the mean lateral-error
RMS by 0.00334~m, the heading-error RMS by 0.00167~rad, and
$J_{\Delta u}$ by 0.12960. Paired intervals supported the heading-error
reduction (95\% CI [0.00026, 0.00334]) and improved control-input
regularity (95\% CI [0.02163, 0.34511]), while no Fusion safety failure
was recorded. These results support uncertainty-adaptive fusion as a
bounded refinement of the task-specific estimator under the matched
simulation protocol.

The six-route grade-necessity experiment evaluated road-grade scheduling
against LPV-MPC path-tracking formulations that do not use grade as a
scheduling variable
\citep{Alcala2020AutonomousRacingLPVMPC,
Tian2022GainScheduledLPVMPC,Cheng2021VelocityVaryingMPC}.
The conditioned grade signal was applied consistently to the prediction
model, measured-disturbance channel, grade feedforward, scheduled
weights, and constraints. Relative to zero-grade LPV-MPC, the
analysis-only truth-driven reference reduced lateral-error RMS by
1.46798~m (95\% CI [0.11715, 3.74682]), heading-error RMS by
0.33065~rad (95\% CI [0.05065, 0.80004]), and the composite
input-increment index by 695.587 (95\% CI [116.408, 1339.454]).
All three route-paired intervals favored the truth-driven reference.
These results demonstrate the control relevance of accurate and
internally consistent grade scheduling under the tested simulation
protocol, while preserving the interpretability and constraint
handling of the model-based LPV-MPC formulation.

Two limitations should be noted. All evidence comes from closed-loop simulation under a fixed evaluation protocol; no physical-vehicle or hardware-in-the-loop test was performed. In addition, no formal stability guarantee is provided for the grade-scheduled closed loop. Future work will focus on three directions. The first is physical validation of the perception-to-control chain on a dual-steering-wheel AGV platform. The second is extending the evaluation to a broader range of routes, load conditions, and measurement disturbances, including confirmation of the lateral-error benefit of the fusion interface. The third is stability and error-bound analysis for LPV-MPC scheduled by learned, uncertainty-adaptive grade estimates.

\FloatBarrier
\balance

\bibliographystyle{IEEEtran}
\bibliography{paper_v1_references}

\EOD

\end{document}
